\documentclass[11pt,a4paper]{amsart}
\usepackage[utf8]{inputenc}
\usepackage[T1]{fontenc}
\usepackage[french]{babel}
\usepackage{amsmath,amssymb,amsthm,amscd}
\usepackage{mathrsfs}
\usepackage{array}
\usepackage{enumerate}
\usepackage{textcomp}
\usepackage[all]{xy}

\theoremstyle{plain}
\newtheorem{theorem}{Theorem}[section]
\newtheorem{proposition}[theorem]{Proposition}
\newtheorem{lemma}[theorem]{Lemme}
\newtheorem{corollary}[theorem]{Corollaire}
\newtheorem{scholie}[theorem]{Scholie}

\theoremstyle{definition}
\newtheorem{definition}[theorem]{Definition}
\newtheorem{remark}[theorem]{Remarque}
\newtheorem{terminology}[theorem]{Terminologie}
\newtheorem{variant}[theorem]{Variante}
\newtheorem{variants}[theorem]{Variantes}
\newtheorem{construction}[theorem]{Construction}
\newtheorem{notation}[theorem]{Notation}

\DeclareMathOperator{\Ind}{Ind}
\DeclareMathOperator{\Res}{Res}
\DeclareMathOperator{\Hom}{Hom}
\DeclareMathOperator{\Ker}{Ker}
\DeclareMathOperator{\Gal}{Gal}
\DeclareMathOperator{\Spec}{Spec}
\DeclareMathOperator{\End}{End}
\DeclareMathOperator{\Aut}{Aut}
\DeclareMathOperator{\Tr}{Tr}
\DeclareMathOperator{\supp}{supp}
\DeclareMathOperator{\Sw}{Sw}
\DeclareMathOperator{\id}{id}
\DeclareMathOperator{\Ima}{Im}

\newcommand{\dbZ}{\mathbb{Z}}
\newcommand{\dbQ}{\mathbb{Q}}
\newcommand{\dbR}{\mathbb{R}}
\newcommand{\dbC}{\mathbb{C}}
\newcommand{\dbA}{\mathbb{A}}
\newcommand{\dbF}{\mathbb{F}}
\newcommand{\dbP}{\mathbb{P}}
\newcommand{\GL}{\mathrm{GL}}
\newcommand{\SL}{\mathrm{SL}}
\newcommand{\PGL}{\mathrm{PGL}}
\newcommand{\dbN}{\mathbb{N}}
\newcommand{\ab}{\mathrm{ab}}
\newcommand{\nr}{\mathrm{nr}}
\DeclareMathOperator{\Inf}{Inf}
\DeclareMathOperator{\sw}{sw}
\DeclareMathOperator{\Ree}{Re}
\def\dim{\mathop{\mathrm{dim}}\nolimits}
\newenvironment{para}{\par\noindent}{\par}
\newenvironment{reduction}{\par\noindent\textbf{R\'eduction. }}{\par}
\newenvironment{variante}{\begin{variant}}{\end{variant}}
\newcommand{\scrR}{\mathscr{R}}
\newcommand{\scrO}{\mathscr{O}}
\newcommand{\scrG}{\mathscr{G}}
\newcommand{\scrF}{\mathscr{F}}
\newcommand{\scrH}{\mathscr{H}}
\newcommand{\scrI}{\mathscr{I}}
\newcommand{\scrJ}{\mathscr{J}}
\newcommand{\scrK}{\mathscr{K}}
\newcommand{\scrL}{\mathscr{L}}
\newcommand{\scrN}{\mathscr{N}}
\newcommand{\scrS}{\mathscr{S}}
\newcommand{\scrV}{\mathscr{V}}
\newcommand{\scrW}{\mathscr{W}}
\newcommand{\scrX}{\mathscr{X}}
\newcommand{\scrY}{\mathscr{Y}}
\newcommand{\scrZ}{\mathscr{Z}}
\newcommand{\calC}{\mathcal{C}}
\newcommand{\calD}{\mathcal{D}}
\newcommand{\calF}{\mathcal{F}}
\newcommand{\calG}{\mathcal{G}}
\newcommand{\calH}{\mathcal{H}}
\newcommand{\calI}{\mathcal{I}}
\newcommand{\calO}{\mathcal{O}}
\newcommand{\calS}{\mathcal{S}}
\newcommand{\calV}{\mathcal{V}}
\newcommand{\calW}{\mathcal{W}}
\newcommand{\fht}{\mathfrak{t}}
\newcommand{\fhp}{\mathfrak{p}}

\newcommand{\ov}[1]{\overline{#1}}
\newcommand{\eps}{\varepsilon}
\newcommand{\fais}{\mathscr{F}}
\newcommand{\wt}{\widetilde}
\newcommand{\wh}{\widehat}
\newcommand{\ti}{\times}
\newcommand{\isom}{\xrightarrow{\sim}}
\newcommand{\inj}{\hookrightarrow}
\newcommand{\surj}{\twoheadrightarrow}
\newcommand{\vl}{\Vert}
\newcommand{\vr}{\Vert}

\title{Les constantes des \'equations fonctionnelles des fonctions $L$}
\author{P. Deligne}
\date{International Summer School on Modular Functions, Antwerp 1972}

\begin{document}
\maketitle
\tableofcontents

\section*{Introduction}
\addcontentsline{toc}{section}{Introduction}
\markboth{INTRODUCTION}{INTRODUCTION}

Les r\'esultats essentiels de cet article sont les suivants.

\medskip
\noindent
\textbf{(A)} On donne une d\'emonstration simple du th\'eor\`eme de Langlands \cite{9} (d\'ej\`a prouv\'e au signe pr\`es par Dwork \cite{6}) qui permet d'\'ecrire la constante des \'equations fonctionnelles des fonctions $L$ d'Artin (ou de Weil \cite{19}) comme produit de constantes locales.

Le th\'eor\`eme \`a d\'emontrer est local, et la d\'emonstration de Langlands avait l'\'el\'egance d'\^etre purement locale. Pour l'essentiel, il s'agit de d\'emontrer une identit\'e (multiplicative) entre sommes de Gauss un peu g\'en\'eralis\'ees chaque fois qu'on a une relation $\dbQ$-lin\'eaire entre caract\`eres de groupes de Galois locaux induits par des caract\`eres de repr\'esentations de dimension un de sous-groupes. Langlands construisait un ensemble g\'en\'erateur de telles relations, et ramenait les identit\'es \`a d\'emontrer \`a des identit\'es entre sommes de Gauss qu'il d\'emontrait en s'appuyant sur le th\'eor\`eme de Stickelberger.

La d\'emonstration pr\'esent\'ee ici est tr\`es proche de la d\'emonstration \cite{5} de l'identit\'e de Hasse-Davenport (cf. aussi Weil \cite{18}). Elle utilise un argument global (4.12), et le fait que les constantes locales \`a d\'efinir sont de nature essentiellement triviale \`a l'infini, dans le cas non ramifi\'e, et dans certains cas tr\`es ramifi\'es (4.14).

\medskip
\noindent
\textbf{(B)} Soient $K$ un corps de fonctions et $\rho: \Gal(\ov{K}/K) \to \GL(n,\dbQ_{\ell})$ une repr\'esentation $\ell$-adique du groupe de Galois de $K$, presque partout non ramifi\'ee. Grothendieck a montr\'e que le produit eul\'erien \'etendu aux places de $K$ (avec $p^{-s}$ remplac\'e par une ind\'etermin\'ee $t$)
\[
Z(\rho,t) = \prod_{v} \det(1 - F_{v} t^{\deg(v)}, \rho(I_{v}))^{-1} \in \dbQ_{\ell}((t))
\]
est une fraction rationnelle, et qu'une \'equation fonctionnelle relie $Z(\rho,t)$ \`a $Z(\check{\rho},t)$:
\[
Z(\check{\rho}, p^{-1}t^{-1}) = \eps \cdot Z(\rho,t)
\]
(la constante $\eps$ est un mon\^ome $a \cdot t^{b}$). Cette th\'eorie est r\'esum\'ee au \S10. Le r\'esultat (A) sugg\`ere une formule pour exprimer la constante $\eps$ comme produit de constantes locales. Nous montrons que cette formule est vraie lorsque $\rho$ appartient \`a un syst\`eme compatible infini de repr\'esentations $\ell$-adiques (9.3 et 9.9). D'apr\`es Weil \cite{16} et Jacquet-Langlands \cite{8}, cette formule, appliqu\'ee au syst\`eme compatible de repr\'esentations $\ell$-adiques d\'efini par une courbe elliptique sur $K$, permet d'associer une forme modulaire pour $\GL(2,K)$ \`a toute courbe elliptique sur $K$.

Il serait tr\`es int\'eressant de savoir si la formule obtenue reste valable pour une quelconque repr\'esentation $\ell$-adique.

\medskip
\noindent
\textbf{Voici le contenu des divers \S\S, et leur interd\'ependance.}

\medskip
Le d\'ebut du \S1 rassemble quelques r\'esultats sur les repr\'esentations induites des groupes finis qui nous seront utiles.

La fin du \S1 (1.11--fin) ne sert pas dans le reste de l'article. Pour $G$ un groupe r\'esoluble, on y d\'ecrit un ensemble $R$ de relations lin\'eaires (sur $\dbZ$) entre les caract\`eres de $G$ induits par des caract\`eres de dimension un de sous-groupes, et on montre que toute telle relation est combinaison lin\'eaire d'\'el\'ements de $R$. Le r\'esultat, et sa d\'emonstration, sont contenus implicitement dans \cite{9}.

Les \S\S2 et 3 servent essentiellement \`a fixer les notations. Au \S2, on rappelle la structure du groupe de Galois d'un corps local \cite{10}, on d\'efinit les ``groupes de Weil'' (variante des groupes de Galois) et on donne quelques \'enonc\'es de la th\'eorie du corps de classe. Le lecteur prendra garde que, dans tout cet article, l'isomorphisme de la th\'eorie du corps de classe local transforme uniformisante en inverse de la substitution de Frobenius (la raison est donn\'ee en 3.6). Le \S3 est un rappel de la th\'ese de Tate, et de Artin \cite{1}.

Le \S4 contient la d\'emonstration du r\'esultat annonc\'e en (A) ci-dessus; le formulaire du \S5 donne une liste de propri\'et\'es fondamentales des constantes locales construites au \S4, et quelques identit\'es entre sommes de Gauss qui en r\'esultent.

Les \S\S6, 7 et 8 (qui ne d\'epend que du \S2) sont pr\'eliminaires \`a la d\'emonstration de (B), donn\'ee au \S9. Dans les \'enonc\'es des \S\S3 et 4, nous avions \'et\'e tr\`es attentifs \`a n'utiliser que des formules ``rationnelles'', ne contenant ni passage \`a la limite ni racine carr\'ee (f\^ut-ce d'un entier positif). Nous en r\'ecoltons le fruit aux \S\S6 et 7. Au \S6, pour $K$ un corps local non archim\'edien et $V$ une repr\'esentation modulaire de $W(\ov{K}/K)$ sur un corps $\Lambda$ de caract\'eristique $\ell \not= p$, nous d\'efinissons (par transport de structure pour $\ell = 0$ et par r\'eduction mod $\ell$ pour $\ell \not= 0$) une constante locale $\eps_{0}(V,\psi,dx) \in \Lambda^{*}$.

Au \S7, pour $K$ global de caract\'eristique $p$ et $V$ une repr\'esentation modulaire de $W(\ov{K}/K)$, nous prouvons par r\'eduction mod $\ell$ une \'equation fonctionnelle pour la fonction $L \in \Lambda(t)$ correspondante.

Au \S8, nous donnons des classes d'isomorphie de repr\'esentations $\ell$-adiques du groupe de Weil d'un corps local une description purement alg\'ebrique, ind\'ependante de la topologie de $\dbQ_{\ell}$. Ceci sert \`a d\'efinir la compatibilit\'e de repr\'esentation $\ell$ et $\ell'$-adique. La notion introduite est plus fine que celle de Serre.

Le \S9 donne la d\'emonstration de (B) (9.3 et 9.9). L'id\'ee est que pour d\'emontrer une identit\'e, il suffit de la d\'emontrer mod $\ell$ pour une infinit\'e de $\ell$. Nous donnons aussi, pour les corps de fonctions, une r\'eponse partielle \`a une question de Serre en montrant (9.8) que deux repr\'esentations $\ell$ et $\ell'$-adiques, donnant lieu au m\^eme polyn\^ome caract\'eristique de Frobenius (suppos\'e dans $\dbQ[t]$) en presque toutes les places, v\'erifient une compatibilit\'e aux places r\'esiduelles.

Le \S10 r\'esume la th\'eorie de Grothendieck des fonctions $L$ (cf. \cite{7} et SGA 5), qui nous a servi de fa\c{c}on cruciale au \S9. En 10.11, je tente d'expliquer le sel de sa m\'ethode. Les r\'esultats farfelus 10.n r\'esultent aussit\^ot de ce qui pr\'ec\`ede. J'ignore leur signification.


\section{Repr\'esentations virtuelles d'un groupe fini}

\begin{para}
Soit $K$ un corps (commutatif). Le cas le plus important pour nous sera celui o\`u $K$ est alg\'ebriquement clos de caract\'eristique $0$, par exemple $K = \dbC$. Soit $G$ un groupe. On notera $R_{K}(G)$ le groupe de Grothendieck de la cat\'egorie des $K[G]$-modules de dimension finie sur $K$. Il s'identifie au $\dbZ$-module libre de base l'ensemble des classes d'isomorphie de repr\'esentations lin\'eaires $K$-irrr\'eductibles de $G$ sur $K$. C'est un anneau commutatif \`a unit\'e (pour le produit tensoriel sur $K$) et m\^eme un $\lambda$-anneau au sens de Grothendieck. Ses \'el\'ements sont les \textbf{repr\'esentations virtuelles} de $G$ sur $K$.
\end{para}

\begin{para}
Pour d\'efinir un homomorphisme $f$ de $R_{K}(G)$ dans un groupe commutatif $X$, il suffit:
\begin{enumerate}[(a)]
\item de d\'efinir $f(V)$, pour $V$ une repr\'esentation de $G$ sur $K$;
\item de v\'erifier que pour toute suite exacte de repr\'esentations $0 \to V' \to V \to V'' \to 0$, on a $f(V) = f(V') + f(V'')$.
\end{enumerate}
\end{para}

\begin{para}
Appliquant cette r\`egle, on d\'efinit la dimension d'une repr\'esentation virtuelle
\[
\dim: R_{K}(G) \longrightarrow \dbZ
\]
et son d\'eterminant
\[
\det: R_{K}(G) \longrightarrow \Hom(G, K^{*}).
\]
Si $H$ est un sous-groupe d'indice fini de $G$, on d\'efinit l'induction comme correspondant \`a l'extension des scalaires $V \mapsto K[G] \otimes_{K[H]} V$ (c'est un foncteur exact; il passe donc aux repr\'esentations virtuelles). Pour $f: G \to H$, on d\'efinit la restriction
\[
\Res^{G}_{H}: R_{K}(G) \longrightarrow R_{K}(H)
\]
comme restriction des scalaires de $K[H]$ \`a $K[G]$. Pour $f$ surjectif, on parle parfois plut\^ot d'inflation.
\end{para}

\begin{para}
Pour tout groupe $G$ on note $G^{\ab}$ le plus grand quotient ab\'elien de $G$. Un $K$-quasi-caract\`ere $\chi$ de $G$ est un homomorphisme $\chi: G \to K^{*}$; on l'identifie au $K$-quasi-caract\`ere de $G^{\ab}$ par lequel il se factorise, et on note $[\chi]$ la repr\'esentation de dimension $1$ correspondante.
\end{para}

\begin{para}
Pour $G$ fini, nous emploierons indiff\'eremment les mots ``caract\`ere'' et ``quasi-caract\`ere''. ``Caract\`ere'' ne signifiera jamais ``caract\`ere d'une repr\'esentation de dimension $> 1$''.
\end{para}

\begin{proposition}[P. Gallagher \cite{20}]
\label{prop1.2}
Soit $G$ un groupe, $H$ un sous-groupe d'indice fini et $t: G^{\ab} \to H^{\ab}$ le transfert. Pour $\rho$ une repr\'esentation virtuelle de dimension $0$ de $H$ et $x \in G$, on a
\[
\det(\Ind^{G}_{H}(\rho))(x) = \det(\rho)(t(x)).
\]
\end{proposition}

On montrera (\`a peine) plus g\'en\'eralement que, si $\theta$ est le d\'eterminant de la repr\'esentation de permutation de $G$ sur $G/H$ ($\theta: G \to \{\pm 1\}$), on a, pour $\rho$ de dimension quelconque,
\begin{equation}\label{eq1.2}
\det(\Ind^{G}_{H}(\rho))(x) = \theta(x)^{\dim \rho} \det(\rho)(t(x)).
\end{equation}

C'est assez \'evident en termes topologiques. Soient $B_{G}$ le classifiant de $G$ et $B_{H}$ celui de $H$, vu comme rev\^etement \`a $[G:H]$ feuillets de $B_{G}$: $f: B_{H} \to B_{G}$. Les repr\'esentations de $G$ s'identifient aux syst\`emes locaux de $K$-vectoriels de dimension finie sur $B_{G}$, et $\Ind^{G}_{H}$ au foncteur $f_{*}$. Le transfert $H^{1}(B_{G},\dbZ) \to H^{1}(B_{H},\dbZ)$ est transpose aux morphismes trace $f_{!}: H^{1}(B_{G},X) \to H^{1}(B_{H},X)$, $a \mapsto \prod_{i=1}^{[G:H]} x_{i}^{*}a$ ($X$ groupe ab\'elien). Pour $X = K^{*}$, $f_{!}$ correspond \`a $\mathbf{N}$: si $\eta \in H^{1}(B_{G},K^{*})$ est la classe d'un syst\`eme local de rang $1$ de $K$-vectoriels $L$, $f_{!}\eta$ est la classe de $\mathbf{N}(L)$, syst\`eme local sur $B_{G}$ qui localement sur $B_{G}$ est le produit tensoriel des images r\'eciproques de $L$ par $[G:H]$ sections disjointes.

\begin{para}
Posons $\det V = \bigwedge^{\dim V} V$ (pour $V$ un syst\`eme local) et soit $\theta = \det f_{*}\dbZ$ le syst\`eme local d\'eterminant de $f$. La proposition \ref{prop1.2} r\'esulte de l'existence d'un isomorphisme canonique, pour $V$ syst\`eme local sur $B_{H}$,
\[
\det f_{*}V \cong \theta^{\otimes \dim V} \otimes \mathbf{N}(\det V).
\]
La construction de cet isomorphisme est locale sur $B_{G}$; elle se ram\`ene \`a la
\end{para}

\begin{construction}
Pour $(V_{i})_{i \in I}$ une famille finie de $K$-vectoriels de dimension $d$, et $e = \det(K^{I})$, on a canoniquement
\[
\det(\textstyle\bigotimes_{i} V_{i}) \cong e^{\otimes d} \otimes \textstyle\bigotimes_{i} \det(V_{i}).
\]
La fl\`eche est, pour $i_{1} \ldots i_{n}$ la suite des \'el\'ements de $I$,
\[
\textstyle\bigwedge_{i}(v_{i1} \wedge \cdots \wedge v_{id}) \longmapsto (e_{i_{1}} \wedge \cdots \wedge e_{i_{n}})^{\otimes d} \otimes \textstyle\bigotimes_{i}(v_{i1} \wedge \cdots \wedge v_{id}).
\]
\end{construction}

\begin{para}
Plut\^ot qu'un classifiant $B_{G}$, on e\^ut plus intrins\`equement pu prendre dans le raisonnement pr\'ec\'edent le topos classifiant de Grothendieck (SGA 4 IV 2.4 p.~315 (pour $E = (\text{Ens})$)).
\end{para}

\begin{para}
Supposons $G$ fini. Nous dirons que $K$ est assez gros s'il contient toutes les racines de l'unit\'e d'ordre divisant le ppcm des ordres des \'el\'ements de $G$. Un groupe $H$ sera dit $p$-\'el\'ementaire s'il est produit d'un $p$-groupe par un groupe cyclique; \'el\'ementaire s'il est $p$-\'el\'ementaire pour un nombre premier $p$ convenable. Nous aurons besoin de la variante suivante du th\'eor\`eme de Brauer sur les caract\`eres induits (qu'on obtient en omettant partout ``de dimension $0$''):
\end{para}

\begin{proposition}\label{prop1.5}
Si $K$ est assez gros, toute repr\'esentation virtuelle de dimension $0$ de $G$ est somme de repr\'esentations virtuelles $\Ind_{H}^{G}(\chi)$, avec $H$ \'el\'ementaire et $\chi$ virtuel de dimension $0$, $\chi$ combinaison lin\'eaire de repr\'esentations de dimension $1$ de $H$.
\end{proposition}

\begin{proof}
D'apr\`es le th\'eor\`eme de Brauer appliqu\'e \`a la repr\'esentation triviale $1_{G}$ de $G$, il existe une d\'ecomposition $1_{G} = \sum_{H} \Ind_{H}^{G}(z_{H})$ ($H$ \'el\'ementaire). Pour $y$ repr\'esentation virtuelle de dimension $0$ de $G$, on a alors $y = \sum_{H} \Ind_{H}^{G}(\Res_{H}^{G}(y) \cdot z_{H})$ avec $\dim(\Res_{H}^{G}(y) \cdot z_{H}) = 0$. Ceci nous ram\`ene \`a supposer $G$ \'el\'ementaire.

Par lin\'earit\'e, on peut supposer $y$ de la forme $\rho - \dim(\rho) \cdot 1_{G}$, avec $\rho$ irr\'eductible, donc (puisque $G$ est nilpotent) induite par un caract\`ere $\chi$ d'un sous-groupe $H$ de $G$ contenant le centre $Z$. On a alors
\[
y = \Ind_{H}^{G}([\chi] - 1_{H}) + (\Ind_{H}^{G}(1_{H}) - [G:H] \cdot 1_{G}).
\]
Le premier terme est du type voulu, et le second est l'inflation d'une repr\'esentation virtuelle de dimension $0$ de $G/Z$. On conclut par une r\'ecurrence sur l'ordre de $G$ (noter que $|G/Z| < |G|$ si $G \not= \{e\}$).
\end{proof}

\begin{scholie}\label{scholie1.6}
Un homomorphisme $f: R_{K}(G) \to X$ est connu quand on conna\^it sa valeur en $1_{G}$ et en les $\Ind_{H}^{G}(\chi - 1_{H})$, pour $H$ \'el\'ementaire et $\chi$ un caract\`ere de $H$.
\end{scholie}

\begin{para}
Soit $A$ un anneau de valuation discr\`ete d'in\'egale caract\'eristique $p$, de corps des fractions $K$ et de corps r\'esiduel $k = A/\mathfrak{m}$. Rappelons la d\'efinition de l'homomorphisme de d\'ecomposition $d: R_{K}(G) \to R_{k}(G)$ (\cite{11} III 2.2): pour $V$ une repr\'esentation de $G$ sur $K$ et $E$ un r\'eseau $G$-invariant, $d([V]) = [E \otimes_{A} k]$.
\end{para}

\begin{proposition}\label{prop1.8}
On suppose $K$ assez gros. Le noyau de l'homomorphisme de d\'ecomposition $d$ est alors engendr\'e par les
\[
\Ind_{H}^{G}([\chi'] - [\chi'']),
\]
pour $H$ un sous-groupe \'el\'ementaire et $\chi', \chi'' \in \Hom(H,K^{*})$ congrus mod $\mathfrak{m}$.
\end{proposition}

\begin{proof}
L'homomorphisme d'anneaux $d$ commutant \`a l'induction, le raisonnement de \ref{prop1.5} nous ram\`ene au cas o\`u $G$ est \'el\'ementaire, donc produit d'un $p$-groupe $P$ par un groupe $H$ d'ordre premier \`a $p$. Les corps $K$ et $k$ \'etant assez gros, on a $R_{K}(G) = R_{K}(P) \otimes R_{K}(H)$, et de m\^eme pour $R_{k}$. Pour un groupe d'ordre premier \`a $p$, $d$ est bijectif. On a donc
\[
R_{K}(G)/\Ker(d) \isom R_{k}(G)
\]
et $\Ker(d)$ est engendr\'e par $\Ker(R_{K}(P) \to R_{k}(P))$ tensoris\'e par $R_{K}(H)$.

Appliquant le th\'eor\`eme de Brauer \`a $H$, on voit qu'il suffit de prouver \ref{prop1.8} pour un $p$-groupe. Ce cas est trivial, car tout caract\`ere $\chi \in \Hom(P,K^{*})$ d'un $p$-groupe est congru \`a $1$ mod $\mathfrak{m}$.
\end{proof}

\begin{para}
Dans la fin de ce num\'ero, on suppose que $K = \dbC$, on ne consid\`ere que des groupes finis et on abr\`eve $R_{K}(G)$ en $R(G)$. Pour $G$ commutatif, on note $G^{*}$ le groupe des caract\`eres de $G$.
\end{para}

\begin{para}
Soit $R_{+}(G)$ le $\dbZ$-module libre de base l'ensemble des classes de $G$-conjugaison de couples $(H,\chi)$: $H$ sous-groupe de $G$ et $\chi$ caract\`ere de $H$. On d\'efinit dans $R_{+}(G)$ une multiplication par la r\`egle
\begin{equation}\label{eq1.9.1}
(A,\alpha) \cdot (B,\beta) = \sum_{(x,y) \in G \setminus (G/A \times G/B)} (A^{x} \cap B^{y}, (\alpha^{x}|_{A^{x} \cap B^{y}}) \cdot (\beta^{y}|_{A^{x} \cap B^{y}}))
\end{equation}
($\langle x,y \rangle$ parcourt des repr\'esentants des orbites de $G$ dans $G/A \times G/B$, $A^{x} = xAx^{-1}$ et $\alpha^{x}$ est le caract\`ere $a \mapsto \alpha(x^{-1}ax)$ de $A^{x}$). Pour cette multiplication, $R_{+}(G)$ est un anneau commutatif d'unit\'e $(G,1)$, et on d\'efinit un homomorphisme d'anneaux
\begin{equation}\label{eq1.9.2}
b: R_{+}(G) \longrightarrow R(G)
\end{equation}
par $b((H,\chi)) = \Ind_{H}^{G}(\chi)$. Cet homomorphisme est surjectif d'apr\`es le th\'eor\`eme de Brauer. Pour $\chi$ un caract\`ere de $G$, la multiplication par $(G,\chi)$ se note $\cdot \chi$ et s'appelle \textbf{torsion} par $\chi$:
\begin{equation}\label{eq1.9.3}
(H,\varphi) \cdot \chi = (H, \varphi \cdot (\chi|_{H})).
\end{equation}
\end{para}

\begin{para}
Pour $G' \subset G$, on d\'efinit l'induction
\begin{equation}\label{eq1.9.4}
\Ind_{G'}^{G}: R_{+}(G') \longrightarrow R_{+}(G)
\end{equation}
par $\Ind_{G'}^{G}((H,\chi)) = (H,\chi)$ (o\`u dans le premier membre $H \subset G'$ et dans le second $H \subset G$). Pour $u: G' \to G$ un morphisme, on d\'efinit la restriction
\begin{equation}\label{eq1.9.5}
\Res_{G'}^{G}: R_{+}(G) \longrightarrow R_{+}(G')
\end{equation}
par $\Res_{G'}^{G}((H,\chi)) = \sum_{s \in u(G') \setminus G/H} (u^{-1}(sHs^{-1}), \chi \circ u)$ ($s$ parcourt un ensemble de repr\'esentants des doubles classes). On dispose de diagrammes commutatifs (pour le second, cf.~\cite{11} p.~75)
\[
\begin{CD}
R_{+}(G') @>{\Ind_{G'}^{G}}>> R_{+}(G) \\
@V{b}VV @VV{b}V \\
R(G') @>{\Ind_{G'}^{G}}>> R(G)
\end{CD}
\qquad
\begin{CD}
R_{+}(G') @<{\Res_{G'}^{G}}<< R_{+}(G) \\
@V{b}VV @VV{b}V \\
R(G') @<{\Res_{G'}^{G}}<< R(G)
\end{CD}
\]
et on a la formule, pour $R \in R_{+}(G)$,
\begin{equation}\label{eq1.9.6}
R \cdot (H,\chi) = \Res_{H}^{G}(R) \cdot \chi \quad \text{dans } R_{+}(H).
\end{equation}
On parle d'inflation lorsque $u$ est surjectif:
\begin{equation}\label{eq1.9.7}
\Inf: R_{+}(G) \longrightarrow R_{+}(G').
\end{equation}
\end{para}

\begin{notation}
Pour toute fonction additive de repr\'esentation $F$, on note de m\^eme la fonction qu'elle d\'efinit sur $R(G)$, et le compos\'e $F \circ b$. Ainsi, $\dim(V) = \dim(b(V))$, \dots
\end{notation}

\begin{variants}
\ \begin{enumerate}[(i)]
\item On note $R_{+}^{0}(G)$ le sous-groupe de $R_{+}(G)$ engendr\'e par les $(H,\chi) - (H,1)$. Pour $(H,\chi)$ parcourant la base canonique de $R_{+}(G)$, \`a l'exception des $(H,1)$, ces \'el\'ements forment une base de $R_{+}^{0}(G)$. On v\'erifie que $R_{+}^{0}(G)$ est un id\'eal de $R_{+}(G)$. Soit $R^{0}(G) \subset R(G)$ l'ensemble des repr\'esentations virtuelles de dimension $0$ de $G$. D'apr\`es \ref{prop1.5}, l'application naturelle, induite par (\ref{eq1.9.2}), de $R_{+}^{0}(G)$ dans $R^{0}(G)$ est surjective.
\item Pour $G$ un groupe topologique (fini ou non), on note $R_{+}(G)$ et $R_{+}^{0}(G)$ les groupes analogues aux pr\'ec\'edents obtenus en se limitant aux sous-groupes ferm\'es $H$ d'indice fini et aux quasi-caract\`eres continus. Les fonctorialit\'es \ref{eq1.9.4}--\ref{eq1.9.7} se g\'en\'eralisent, mutatis mutandis.
\end{enumerate}
\end{variants}

\begin{para}
La fin de ce \S\ ne servira plus par la suite. Nous y exposons des r\'esultats tir\'es de Langlands \cite{9}.
\end{para}

\begin{lemma}\label{lem1.11}
Supposons que $G = H \cdot C$, avec $C$ commutatif distingu\'e, et soit $\chi$ un caract\`ere de $H$. Pour tout $\mu \in C^{*}$, soit $G_{\mu} \subset G$ le fixateur de $\mu$ (dans $C^{*}$ par conjugaison). Si $\chi|_{H \cap C} = \mu|_{H \cap C}$, on note $\{\chi,\mu\}$ le caract\`ere de $G_{\mu} = (H \cap G_{\mu}) \cdot C$ qui prolonge $\chi|_{H \cap G_{\mu}}$ et $\mu$. On a, pour $\mu$ parcourant un ensemble de repr\'esentants des classes de $G$-conjugaison de caract\`eres de $C$ tels que $\chi|_{H \cap C} = \mu|_{H \cap C}$,
\[
\Ind_{H}^{G}(\chi) = \sum_{\mu \in C^{*}/G} \Ind_{G_{\mu}}^{G}(\{\chi,\mu\}).
\]
La v\'erification est laiss\'ee au lecteur.
\end{lemma}

\begin{para}
Un \'el\'ement $R$ de $\Ker(R_{+}(G) \xrightarrow{b} R(G))$ sera dit de type I, II ou III s'il existe un quotient $A$ d'un sous-groupe $B$ de $G$ tel que $R$ s'obtienne \`a partir d'un \'el\'ement $S$ de l'un des types respectifs suivants de $\Ker(R_{+}(A) \to R(A))$ par inflation de $A$ \`a $B$, torsion par un caract\`ere de $B$ et induction de $B$ \`a $G$. On d\'esigne par $\ell$ un nombre premier.
\end{para}

\begin{para}\ \begin{enumerate}
\item[\textbf{I}] $A = \dbZ/\ell\dbZ$;
\[
S = (e,1) - \sum_{\chi \in A^{*}} (A,\chi).
\]
\item[\textbf{II}] $A$ est une extension centrale de $(\dbZ/\ell\dbZ)^{2}$ par un groupe ab\'elien $Z$, $H_{i}$ ($i=1,2$) l'image r\'eciproque dans $A$ des premiers et second facteurs $\dbZ/\ell\dbZ$, et $\chi_{i}$ un caract\`ere de $H_{i}$. On suppose que $\chi_{1}|_{Z} = \chi_{2}|_{Z}$, et que ce caract\`ere de $Z$ est non trivial sur un commutateur.
\[
\xymatrix{
& H_{1} \ar@{^{(}->}[dr] & \\
\dbZ/\ell \times \dbZ/\ell & Z \ar@{^{(}->}[ul] \ar@{^{(}->}[u] \ar@{^{(}->}[ur] & A \ar@{->>}[l]
}
\]
La relation $S$ est
\[
S = (H_{1},\chi_{1}) - (H_{2},\chi_{2}) + \cdots
\]
\item[\textbf{III}] $A$ est un produit semi-direct $H \cdot C$, $\chi$ est un caract\`ere de $H$ et le sous-groupe distingu\'e $C$ est minimal parmi les sous-groupes commutatifs distingu\'es non triviaux de $A$. Pour tout caract\`ere $\mu$ de $C$, soit $A_{\mu}$ le fixateur de $\mu$. Notons $\{\chi,\mu\}$ le caract\`ere de $A_{\mu}$ qui prolonge $\chi$ et $\mu$. $S$ exprime la relation suivante, o\`u $\mu$ parcourt un ensemble de repr\'esentants des orbites de $A$ agissant dans $C^{*}$ par conjugaison:
\[
S = \sum_{\mu \in C^{*}/A} \Ind_{A_{\mu}}^{A}(\{\chi,\mu\}) - \cdots
\]
\end{enumerate}
Ces relations sont des cas particuliers de \ref{lem1.11} (pour I, $H = \{e\}$, $C = A$; pour II, $H = H_{1}$, $C = H_{2}$; pour III, $H$ et $C$ sont $H$ et $C$).
\end{para}

\begin{para}
Toute relation d\'eduite par induction, inflation ou torsion d'une relation de type I (resp.~II, III) est encore du m\^eme type.
\end{para}

\begin{theorem}[Langlands \cite{9} \S18]\label{thm1.13}
Si $G$ est nilpotent, le noyau de $b: R_{+}(G) \to R(G)$ est engendr\'e (comme $\dbZ$-module) par les relations de type \textup{I} et \textup{II}.
\end{theorem}

\begin{lemma}\label{lem1.13.1}
Si $G$ est commutatif, $\Ker(b)$ est engendr\'e par les relations de type \textup{I}.
\end{lemma}

\begin{proof}
Il faut montrer que pour tout sous-groupe $H$ de $G$ et tout $\chi \in H^{*}$, la relation
\begin{equation}\label{eq1.13.1}
(G,\chi) - \Ind_{H}^{G}(\chi) = 0
\end{equation}
est cons\'equence ($=$ combinaison lin\'eaire) de relations de type I. Celles-ci sont les suivantes: pour $H' \subset H'' \subset G$ avec $[H'':H']$ premier, et $\chi$ un caract\`ere de $H'$,
\[
\Ind_{H'}^{G}(\chi) = \Ind_{H''}^{G}(\Ind_{H'}^{H''}(\chi)).
\]
Que (\ref{eq1.13.1}) en r\'esulte se voit en prenant une suite de Jordan-H\"older de $G/H$.
\end{proof}

\begin{lemma}\label{lem1.13.2}
Soient $C$ un sous-groupe commutatif distingu\'e de $G$, $T$ un ensemble de repr\'esentants de $C^{*}/G$, $\mu \in T$, et $G_{\mu}$ le fixateur de $\mu$. Soit $(H_{i})$ une famille de sous-groupes de $G$ contenant $C$, et $\chi_{i}$ un caract\`ere de $H_{i}$. Si la relation
\[
\sum_{i} n_{i} \Ind_{H_{i}}^{G}(\chi_{i}) = 0
\]
est vraie dans $R(G)$, la relation
\[
\sum_{\chi_{i}|_{C} = \mu} n_{i} \Ind_{H_{i} \cap G_{\mu}}^{G_{\mu}}(\chi_{i}|_{H_{i} \cap G_{\mu}}) = 0
\]
est vraie dans $R(G_{\mu})$.
\end{lemma}

La v\'erification du lemme suivant est laiss\'ee au lecteur.

\begin{proof}[D\'emonstration du th\'eor\`eme \ref{thm1.13}]
Soit $Z$ le centre de $G$ et prouvons \ref{thm1.13} par r\'ecurrence sur l'ordre de $G/Z$. D'apr\`es \ref{lem1.13.1}, on peut supposer $G$ non commutatif. Soit donc $C$ un sous-groupe commutatif distingu\'e contenant $Z$, tel que $[C:Z] = \ell$ soit premier, et d'image centrale dans $G/Z$.

Soit une relation
\begin{equation}\label{eq1.13.3}
\sum_{i} n_{i} (H_{i},\chi_{i}) = 0.
\end{equation}
Elle est cons\'equence de relations induites par les relations
\[
\sum_{\chi_{i}|_{H_{i} \cap C = \chi_{i}}} [\chi_{i}]
\]
(qui sont combinaisons lin\'eaires de relations de type I: on peut le d\'eduire de \ref{lem1.13.1} appliqu\'e \`a $H_{i}Z/\Ker(\chi_{i})$) et d'une relation (2) analogue \`a (\ref{eq1.13.3}), avec $H_{i} \supset Z$.

Les relations (a$_{i}$) sont combinaisons lin\'eaires de relations de type I. La relation (2) est cons\'equence de relations induites par des relations \ref{lem1.11} pour $HC$, $H$, $C$, $\chi$ (avec $H \supset Z$)
\[
\text{(b)} \qquad \Ind_{H}^{HC}(\chi) = \sum_{\substack{\mu \in C^{*} \\ \chi|_{H \cap C} = \mu|_{H \cap C}}} \Ind_{G_{\mu}}^{HC}(\{\chi,\mu\})
\]
et d'une relation (3) analogue \`a (\ref{eq1.13.3}), avec cette fois $H_{i} \supset C$. Appliquons \ref{lem1.13.2} \`a (3), pour l'exprimer comme somme de relations induites par des relations
\[
\text{(4)} \qquad \sum_{i} n_{i} \Ind_{H_{i} \cap G_{\mu}}^{G_{\mu}}(\chi_{i}|_{H_{i} \cap G_{\mu}}) = 0.
\]
L'image de $C$ dans $G_{\mu}/\Ker(\mu)$ est centrale (car $\mu$, invariant par $G_{\mu}$, est injectif sur $C/\Ker(\mu)$). L'hypoth\`ese de r\'ecurrence s'applique donc \`a $G_{\mu}$, et les relations (4) sont du type voulu. Prouvons que les relations (b) le sont. Si $HC \not= G$, cela r\'esulte de l'hypoth\`ese de r\'ecurrence, appliqu\'ee \`a $HC$. Si $HC = G$, $H$ est distingu\'e, car il contient $Z$ et que $C$ est central mod $Z$. Soit $A$ l'intersection des noyaux des conjugu\'es de $\chi$. La relation (b) est l'inflation d'une relation analogue pour $G/A$. Si l'hypoth\`ese de r\'ecurrence s'applique \`a $G/A$, on a gagn\'e. Sinon, on se retrouve dans la m\^eme situation qu'avant, avec $A = \{e\}$, et on conclut par le
\end{proof}

\begin{lemma}\label{lem1.13.3}
Avec les notations pr\'ec\'edentes, si $A = \{e\}$, la relation (b) est de la forme \textup{II}.
\end{lemma}

\begin{proof}
Le groupe $H$ est commutatif, car des caract\`eres s\'eparent ses \'el\'ements. Le groupe $Z$ est cyclique, car un seul caract\`ere s\'epare ses \'el\'ements. Pour $c \in C$ et $h \in H$, posons $u(c,h) = chc^{-1}h^{-1}$. C'est un \'el\'ement de $Z$, car $C$ est central mod $Z$. Il d\'epend biadditivement de $c$ et $h$ et d\'efinit $u: C/Z \otimes H/Z \to Z$. Soit $c$ engendrant $C/Z$. Puisque $Z$ est le centre et que $G \not= Z$, $u(c,\cdot): H/Z \to Z$ est injectif, d'image cyclique tu\'ee par $\ell$ et non triviale. D\`es lors, $|C/Z| = |H/Z| = \ell$, et $G$ est extension centrale de $C/Z \times H/Z$ par $Z$. Le caract\`ere $\chi$ est distinct d'au moins un de ses conjugu\'es (sinon $H$ serait central); il est donc non trivial sur un commutateur, et on est dans la situation II.
\end{proof}

\begin{theorem}[Langlands \cite{9}]\label{thm1.14}
Si $G$ est r\'esoluble, le noyau de $b: R_{+}(G) \to R(G)$ est engendr\'e par les relations de type \textup{I}, \textup{II} et \textup{III}.
\end{theorem}

\begin{lemma}\label{lem1.14.1}
Si $G$ est r\'esoluble, le sous-$\dbZ$-module $N(G)$ de $R_{+}(G)$ engendr\'e par les relations de type \textup{I}, \textup{II} et \textup{III} est un id\'eal.
\end{lemma}

\begin{proof}[D\'emonstration de \ref{thm1.14} et \ref{lem1.14.1}]
Prouvons \ref{thm1.14} et \ref{lem1.14.1} par une r\'ecurrence simultan\'ee sur l'ordre de $G$.

\medskip
\noindent
\textbf{(A)} \ref{thm1.14} pour les sous-groupes propres de $G$ implique \ref{lem1.14.1} pour $G$.

Soient $R \in R_{+}(G)$ de type I, II ou III, $H$ un sous-groupe et $\chi$ un caract\`ere de $G$. Prouvons que $R \cdot (H,\chi) \in N(G)$. Si $H = G$, cela r\'esulte de ce que $N$ est stable par torsion par un caract\`ere de $G$. Si $H \not= G$, on a (\ref{eq1.9.6}):
\[
R \cdot (H,\chi) = \Ind_{H}^{G}(\Res_{H}^{G}(R) \cdot \chi),
\]
o\`u $\Res_{H}^{G}(R) \in \Ker(R_{+}(H) \longrightarrow R(H))$. Vu l'hypoth\`ese, $\Res_{H}^{G}(R) \in N(H)$ et $R \cdot (H,\chi)$, qui se d\'eduit de $\Res_{H}^{G}(R)$ par torsion et induction, est dans $N(G)$.

\medskip
\noindent
\textbf{(B)} \ref{lem1.14.1} pour $G$ $\Rightarrow$ \ref{thm1.14} pour les groupes d'ordre plus petit impliquent \ref{thm1.14} pour $G$.

Soit $Z$ le centre de $G$. D'apr\`es \ref{thm1.13}, on peut supposer (et on suppose) que $G$ n'est pas nilpotent. Distinguons deux cas.

\medskip
\noindent
\textbf{(B1)} $Z \not= \{e\}$.

D'apr\`es le th\'eor\`eme de Brauer dans $G/Z$, il existe des sous-groupes nilpotents $H_{i} \supset Z$ de $G$ et des caract\`eres $\chi_{i}$ de $H_{i}/Z$ tels que
\[
1_{G/Z} = \sum_{i} n_{i} \Ind_{H_{i}/Z}^{G/Z}(\chi_{i}).
\]
L'hypoth\`ese de r\'ecurrence s'applique \`a $G/Z$, de sorte que
\[
S = (G,1) - \sum_{i} n_{i}(H_{i},\chi_{i}) \in N(G).
\]
Soit $R \in \Ker(b)$. On a (\ref{eq1.9.6}):
\[
R = R \cdot S + \sum_{i} n_{i} \Ind_{H_{i}}^{G}(\Res_{H_{i}}^{G}(R) \cdot \chi_{i}).
\]
Puisque $S \in N(G)$, on a $R \cdot S \in N(G)$. On conclut en notant que, d'apr\`es \ref{thm1.13} ou l'hypoth\`ese, $\Res_{H_{i}}^{G}(R) \in N(H_{i})$.

\medskip
\noindent
\textbf{(B2)} $Z = \{e\}$.

Soit $C$ un sous-groupe commutatif non trivial minimal et, distingu\'e; pour $\mu \in C^{*}$, soit $G_{\mu}$ son fixateur. Proc\'edons comme dans la d\'emonstration de \ref{thm1.13}. On trouve que toute relation $R \in \Ker(b)$ est combinaison lin\'eaire de relations (a) induites par des relations \ref{lem1.11} pour $HC$, $H$, $C$, $\chi$ ($H \supsetneqq C$, $\chi$ caract\`ere de $H$, $H \cap C \subsetneqq C$) et de relations induites par des relations exprimant une identit\'e
\[
\text{(b)} \qquad \Ind_{H_{i}}^{G}(\chi_{i}) = \cdots \quad (\chi_{i}|_{C} = \text{caract\`ere de } C).
\]
Puisque $Z = \{e\}$, $|G_{\mu}/\Ker(\mu)| < |G|$ et l'hypoth\`ese de r\'ecurrence s'applique \`a ces derni\`eres.

Consid\'erons une relation (a). Si $HC \not= G$, on applique l'hypoth\`ese de r\'ecurrence. Si $HC = G$, $H \cap C = \{e\}$; $H$ est distingu\'e car distingu\'e dans $H$ et dans $C$. Par minimalit\'e de $C$, $H \cap C = C$; la relation est du type III.
\end{proof}

\begin{remark}\label{rem1.15}
Il est vraisemblablement possible d'extraire de \cite{9} la description d'un syst\`eme g\'en\'erateur de $\Ker(R_{+}^{0}(G) \longrightarrow R^{0}(G))$, pour des groupes convenables $G$. Je n'ai pas eu le courage de m'y essayer.
\end{remark}


\section{Groupes de Weil}

\subsection{Racines de l'unit\'e}

\begin{para}
Soit $K$ un corps s\'eparablement clos d'exposant caract\'eristique $p$. Pour tout entier $n$, on note $\dbZ/n\dbZ(1)(K)$, ou simplement $\dbZ/n\dbZ(1)$, le groupe des racines $n$-i\`emes de l'unit\'e de $K$. Pour $n|m$, on dispose d'une application de transition $\dbZ/m\dbZ(1) \to \dbZ/n\dbZ(1)$, $x \mapsto x^{m/n}$; on pose
\[
\dbZ(1) = \varprojlim_{n} \dbZ/n\dbZ(1).
\]
Pour $\ell$ un nombre premier, on pose de m\^eme
\[
\dbZ_{\ell}(1) = \varprojlim_{k} \dbZ/\ell^{k}\dbZ(1).
\]
Les applications \'evidentes $\dbZ/n\dbZ(1) \to \dbZ/\ell^{n}\dbZ(1)$ induisent un isomorphisme
\[
\dbZ(1) \otimes_{\dbZ} \dbZ_{\ell} \isom \dbZ_{\ell}(1).
\]
Pour tout $\dbZ$-module $V$, on pose $V(1) = V \otimes_{\dbZ} \dbZ(1)$. Pour $V$ un $\dbZ_{\ell}$-module, on a $V(1) = V \otimes_{\dbZ_{\ell}} \dbZ_{\ell}(1)$. Pour le $\dbZ$-module $\dbQ/\dbZ$, $\dbQ/\dbZ(1)$ est canoniquement isomorphe au groupe de racines de l'unit\'e de $K$: pour $\frac{m}{n} \in \dbQ$ et $x \in \dbZ(1)$, d'image $\ov{x}$ dans $\dbZ/n\dbZ(1)$, l'image de $\frac{m}{n} \otimes x$ dans $\dbQ/\dbZ(1)$ est identifi\'ee \`a $\ov{x}^{m}$.
\end{para}

\begin{para}
Pour $\ell \not= p$, $\dbZ_{\ell}(1)$ est libre de rang $1$ sur $\dbZ_{\ell}$. Pour $i \in \dbZ$, on note $\dbZ_{\ell}(i)$ sa puissance tensorielle $i$-i\`eme (pour $i \leq 0$, $\dbZ_{\ell}(i)$ est le dual de $\dbZ_{\ell}(-i)$). Pour tout $\dbZ_{\ell}$-module $V$, on pose encore $V(i) = V \otimes_{\dbZ_{\ell}} \dbZ_{\ell}(i)$.
\end{para}

\subsection{Notations}

\begin{para}\label{para2.2.1}
Soit $K$ un corps local. On note $\vl x \vr$ la valeur absolue normalis\'ee de $K$; si $dx$ est une mesure de Haar,
\[
\vl a \vr = \frac{d(ax)}{dx}, \quad \text{i.e. } \int f(ax)\,dx = \vl a \vr^{-1} \int f(x)\,dx.
\]
Pour $s \in \dbC$, on note $\omega_{s}$ le quasi-caract\`ere $x \mapsto \vl x \vr^{s}$ de $K^{*}$ dans $\dbC^{*}$.
\end{para}

\begin{para}
Si $K$ est non archim\'edien ($\dbR$ ou $\dbC$), on note $\scrO$ l'anneau de la valuation $v$ de $K$, $\pi$ une uniformisante, $k$ le corps r\'esiduel $\scrO/(\pi)$, $p$ la caract\'eristique de $k$, et on pose $d = [k:\dbF_{p}]$, $q = p^{d} = |k|$. On a $\vl x \vr = q^{-v(x)}$, donc $x \mapsto \vl x \vr$.
\end{para}

\begin{para}
Supposons $K$ non archim\'edien. On d\'esigne par $\ov{K}$ une cl\^oture s\'eparable de $K$ et on note $\ov{\scrO}$ la cl\^oture int\'egrale de $\scrO$ dans $\ov{K}$, $\ov{k}$ le corps r\'esiduel de $\ov{\scrO}$ (une cl\^oture alg\'ebrique de $k$), et $v$ la valuation de $\ov{K}$, \`a valeurs dans $\dbQ$, qui prolonge celle de $K$.
\end{para}

\begin{para}
On dispose d'une filtration en trois crans du groupe de Galois $\Gal(\ov{K}/K)$, de quotients successifs
\[
\begin{CD}
\Gal(\ov{K}/K) @>>> I @>>> P @>>> 1 \\
@VVV @VVV @VVV \\
\Gal(\ov{k}/k) @>{\sim}>> \hat{\dbZ} @. \text{un pro-$p$-groupe} @.
\end{CD}
\]
On l'obtient comme suit (pour les d\'emonstrations, on renvoie \`a \cite{10}).
\begin{enumerate}[(a)]
\item Par transport de structure, $\Gal(\ov{K}/K)$ agit sur $\ov{\scrO}$ et $\ov{k}$. Le groupe d'inertie $I$ est le noyau de la fl\`eche de r\'eduction $\Gal(\ov{K}/K) \to \Gal(\ov{k}/k)$. Le quotient $\Gal(\ov{k}/k)$ est canoniquement isomorphe \`a $\hat{\dbZ}$, de g\'en\'erateur topologique la substitution de Frobenius $\varphi: x \mapsto x^{q}$.
\item Le groupe d'inertie sauvage $P$ est l'unique pro-$p$-groupe de Sylow de $I$. Le morphisme \'equivariant $t: I \to \dbZ(1)(\ov{k})$, de noyau $P$, est caract\'eris\'e par la congruence (modulo l'id\'eal maximal de $\ov{\scrO}$)
\[
\frac{\sigma(x)}{x} \equiv t(\sigma)^{v(x)} \pmod{\mathfrak{m}_{\ov{\scrO}}} \qquad (\sigma \in I,\; x \in \ov{K}^{*}).
\]
\end{enumerate}
\end{para}

\begin{para}
Pour $p$ premier, on note $t_{p}$ l'application compos\'ee $I \xrightarrow{t} \dbZ(1) \to \dbZ_{p}(1)$.
\end{para}

\begin{para}
Soit $k$ un corps fini \`a $q$ \'el\'ements, de cl\^oture alg\'ebrique $\ov{k}$. On note $W(\ov{k}/k)$ le sous-groupe de $\Gal(\ov{k}/k) \cong \hat{\dbZ}$ engendr\'e par $\varphi: x \mapsto x^{q}$. On a canoniquement
\[
\begin{CD}
W(\ov{k}/k) @>{\sim}>> \dbZ \quad (\text{engendr\'e par } \varphi) \\
\cap @VVV \\
\Gal(\ov{k}/k) @>{\sim}>> \hat{\dbZ}
\end{CD}
\]
Le Frobenius g\'eom\'etrique $F \in W(\ov{k}/k)$ est par d\'efinition $\varphi^{-1}$.
\end{para}

\begin{para}
Soient $K$, $\ov{K}$ comme en \ref{para2.2.1}. Le groupe de Weil $W(\ov{K}/K)$, ou groupe de Weil absolu de $K$, est le sous-groupe de $\Gal(\ov{K}/K)$ des $\sigma$ tels que $v(\sigma)$ soit une puissance enti\`ere du Frobenius $\varphi$. On le munit de la topologie induisant la topologie naturelle de $I$, et pour laquelle $I$ est ouvert. Le groupe $\Gal(\ov{K}/K)$ est le compl\'et\'e de $W(\ov{K}/K)$ pour la topologie des sous-groupes ouverts d'indice fini; pour $L$ une extension finie de $K$ dans $\ov{K}$, le sous-groupe correspondant de $W(\ov{K}/K)$ s'identifie \`a $W(\ov{K}/L)$. On appellera parfois Frobenius g\'eom\'etriques les \'el\'ements de $\Gal(\ov{K}/K)$ ou $W(\ov{K}/K)$ d'image $F$.
\end{para}

\begin{para}
Supposons $K$ archim\'edien, de cl\^oture alg\'ebrique $\ov{K}$. On d\'efinit $W(\ov{K}/K)$ comme l'extension suivante de $\Gal(\ov{K}/K)$ par $\ov{K}^{*}$:
\begin{enumerate}
\item[$1)$] $K = \dbR$: $W(\ov{K}/K)$ est engendr\'e par $\ov{K}^{*}$ et un \'el\'ement $F$ avec $F^{2} = -1$ et $F^{-1}zF = \ov{z}$ pour $z \in \ov{K}^{*}$.
\item[$2)$] $K = \dbC$: $W(\ov{K}/K) = \ov{K}^{*}$.
\end{enumerate}
\end{para}

\subsection{Th\'eorie du corps de classe}

\begin{para}
La th\'eorie du corps de classe local fournit un isomorphisme entre $W(\ov{K}/K)^{\ab}$ et $K^{*}$.
\end{para}

Supposons $K$ non archim\'edien. Pour une raison expliqu\'ee en \ref{para3.6}, nous normaliserons alors cet isomorphisme ($=$ choisirons lui ou son inverse) de telle sorte que les Frobenius g\'eom\'etriques correspondent aux uniformisantes.
\[
\begin{CD}
1 @>>> I @>>> W(\ov{K}/K) @>>> W(\ov{k}/k) @>>> 1 \\
@. @VVV @VVV @VVV @. \\
1 @>>> \scrO^{*} @>>> K^{*} @>{v}>> \dbZ @>>> 1
\end{CD}
\]

Une repr\'esentation de $W(\ov{K}/K)$ est non ramifi\'ee (resp.~mod\'er\'ement ramifi\'ee) si elle est triviale sur $I$ (resp.~sur $P$). De m\^eme, un quasi-caract\`ere $\chi$ de $K^{*}$ est non ramifi\'e (resp.~mod\'er\'ement ramifi\'e) s'il est trivial sur $\scrO^{*}$ (resp.~$1+(\pi)$). Pour $\ell$ premier \`a $p$, l'action de $\Gal(\ov{K}/K)$ sur $\dbZ_{\ell}(1)$ est non ramifi\'ee, d\'ecrite par le quasi-caract\`ere $\omega_{1} = \vl x \vr$ de $K^{*}$.

Pour $K$ complexe, l'isomorphisme de la th\'eorie du corps de classe local est celui \'evident sur \ref{para2.2.1}; pour $K$ r\'eel, c'est celui qui rend vrai \ref{para2.3.1}, \ref{para2.3.2} ci-dessous.

\begin{para}
Soit $L \subset \ov{K}$ une extension finie de $K$, de sorte que $W(\ov{K}/L) \subset W(\ov{K}/K)$. Les diagrammes suivants (o\`u $N$ est la norme et $t$ le transfert) sont commutatifs (\cite{10} XIII \S4):
\begin{equation}\label{eq2.3.1}
\begin{CD}
W(\ov{K}/L) @>{\sim}>> W(\ov{K}/L)^{\ab} @>{\sim}>> L^{*} \\
@VVV @. @VV{N_{L/K}}V \\
W(\ov{K}/K) @>{\sim}>> W(\ov{K}/K)^{\ab} @>{\sim}>> K^{*}
\end{CD}
\end{equation}
\begin{equation}\label{eq2.3.2}
\begin{CD}
W(\ov{K}/K) @>{\sim}>> W(\ov{K}/K)^{\ab} @>{\sim}>> K^{*} \\
@AAA @. @VVV \\
W(\ov{K}/L) @>{\sim}>> W(\ov{K}/L)^{\ab} @>{\sim}>> L^{*}
\end{CD}
\end{equation}
\end{para}

\subsection{Groupes de Weil globaux}

\begin{para}
Soient $K$ un corps de fonctions d'une variable sur $\dbF_{p}$ et $k$ son corps des constantes (fermeture alg\'ebrique de $\dbF_{p}$ dans $K$). Soient $\ov{K}$ une cl\^oture s\'eparable de $K$, $\ov{k}$ la cl\^oture alg\'ebrique de $k$ dans $\ov{K}$ et $\Gal^{0}(\ov{K}/K)$ le noyau de l'application de restriction de $\Gal(\ov{K}/K)$ dans $\Gal(\ov{k}/k)$. Le groupe de Weil global (absolu) $W(\ov{K}/K)$ est, en analogie avec \ref{para2.2.1}, d\'efini par le diagramme
\[
\begin{CD}
0 @>>> \Gal^{0}(\ov{K}/K) @>>> \Gal(\ov{K}/K) @>>> \Gal(\ov{k}/k) @>>> 0 \\
@. @| @AAA @AAA @. \\
0 @>>> \Gal^{0}(\ov{K}/K) @>>> W(\ov{K}/K) @>>> W(\ov{k}/k) @>>> 0
\end{CD}
\]
\end{para}

\begin{para}
Soient $v$ une place de $K$ et supposons choisie une place de $\ov{K}$ au-dessus de $v$. On dispose alors d'un diagramme commutatif
\[
\begin{CD}
0 @>>> I_{v} @>>> W(\ov{K}_{v}/K_{v}) @>>> W(\ov{k(v)}/k(v)) @>>> 0 \\
@. @VVV @VVV @VVV @. \\
0 @>>> \Gal^{0}(\ov{K}/K) @>>> W(\ov{K}/K) @>>> W(\ov{k}/k) @>>> 0
\end{CD}
\]
\end{para}

\begin{para}
La th\'eorie du corps de classe global fournit un isomorphisme
\[
W(\ov{K}/K)^{\ab} \isom \dbA_{K}^{*}/K^{*} \quad (\text{groupe des classes d'id\`eles})
\]
rendant commutatifs les diagrammes
\[
\begin{CD}
W(\ov{K}_{v}/K_{v})^{\ab} @>{\sim}>> K_{v}^{*} \\
@VVV @VVV \\
W(\ov{K}/K)^{\ab} @>{\sim}>> \dbA_{K}^{*}/K^{*}
\end{CD}
\]
\end{para}

\begin{para}
Pour la d\'efinition des groupes de Weil globaux dans le cas des corps de nombres on renvoie \`a Weil \cite{19} ou \`a \cite{2} XIV. Nous ne nous en servirons qu'incidemment.
\end{para}


\section{Rappels sur les fonctions $L$}\label{sec3}

\begin{para}
Soit $K$ un corps local et reprenons les notations \ref{para2.2.1}. On d\'esignera par $dx$ une mesure de Haar sur $K$, par $d^{*}x$ une mesure de Haar sur $K^{*}$ et par $\psi$ un caract\`ere additif non trivial de $K$.
\end{para}

\begin{para}
Pour $\chi$ un quasi-caract\`ere (= homomorphisme continu) $\chi: K^{*} \to \dbC^{*}$, on d\'efinit comme suit $L(\chi) \in \dbC \cup \{\infty\}$.
\end{para}

\begin{para}
\textbf{(3.2.1)} $K = \dbR$. On pose $\Gamma_{\dbR}(s) = \pi^{-s/2} \Gamma(s/2)$, et, pour $x$ le plongement de $K$ dans $\dbC$ et $N = 0$ ou $1$,
\[
L(\chi) = \Gamma_{\dbR}(s) \quad \text{pour } \chi = \omega_{s} \text{ ou } \chi = \text{signe} \cdot \omega_{s} \text{ (avec } N = 0 \text{ ou } 1).
\]
\end{para}

\begin{para}
\textbf{(3.2.2)} $K = \dbC$. On pose $\Gamma_{\dbC}(s) = 2(2\pi)^{-s} \Gamma(s)$, et, pour $z$ un plongement de $K$ dans $\dbC$ et $N \geq 0$,
\[
L(\chi) = \Gamma_{\dbC}(s) \quad \text{pour } \chi = z^{-N} \ov{z}^{-M} \omega_{s} \text{ avec } N+M \text{ donn\'e}.
\]
\end{para}

\begin{para}
\textbf{(3.2.3)} $K$ non archim\'edien. On pose
\[
L(\chi) = \begin{cases}
1 & (\chi \text{ ramifi\'e}) \\
\dfrac{1}{1 - \chi(\pi)} & (\chi \text{ non ramifi\'e}).
\end{cases}
\]
\end{para}

\begin{para}
$V$ et $dx$ etant donn\'es, on pose, pour $f$ une fonction $C^{\infty}$ \`a support compact sur $K$,
\[
\wh{f}(y) = \int f(x) \psi(xy)\, dx.
\]
On d\'efinit $\eps(\chi,\psi,dx) \in \dbC^{*}$ par l'\'equation fonctionnelle locale de Tate (ind\'ependante de $f$)
\begin{equation}\label{eq3.3.1}
\frac{\int \wh{f}(x) \chi^{-1}(x) d^{*}x}{L(\chi^{-1})} = \eps(\chi,\psi,dx) \frac{\int f(x) \chi(x) d^{*}x}{L(\chi)}
\end{equation}
(on prend la m\^eme mesure de Haar $d^{*}x$ dans les deux membres). Les deux membres sont d\'efinis par prolongement analytique en $\chi$ dans les familles $\chi \cdot \omega_{s}$ ($s \in \dbC$). Dans ces familles, $\eps$ est de la forme $A \cdot e^{Bs}$. On a
\begin{equation}\label{eq3.3.2}
\eps(\chi,\psi(ax),dx) = \chi(a) \cdot \vl a \vr^{-1} \cdot \eps(\chi,\psi,dx)
\end{equation}
\begin{equation}\label{eq3.3.3}
\eps(\chi,\psi,dx) = \lambda \cdot \eps(\chi,\psi,dx') \quad \text{pour } dx' = \lambda \cdot dx.
\end{equation}
\end{para}

\begin{para}
Pour $K$ non archim\'edien, nous poserons
\begin{itemize}
\item $n(\psi)$ = le plus grand entier $n$ tel que $\psi|_{\pi^{-n}\scrO} = 1$;
\item $a(\chi)$ = le conducteur de $\chi$ ($0$ si $\chi$ est non ramifi\'e, le plus petit entier $m$ tel que $\chi|_{1+(\pi)^{m}} = 1$ si $\chi$ est ramifi\'e);
\item $\sw(\chi) = 0$ pour $\chi$ non ou mod\'er\'ement ramifi\'e, $= a(\chi)-1$ pour $\chi$ ramifi\'e;
\item $y$ = un \'el\'ement de $K^{*}$ de valuation $n(\psi) + \sw(\chi) + 1$.
\end{itemize}
La constante locale $\eps$ est donn\'ee par (\ref{eq3.3.2}), (\ref{eq3.3.3}) et les formules suivantes.
\end{para}

\begin{para}
\textbf{(3.4.1)} $K = \dbR$. $N = 0$ ou $1$. Soient $x$ le plongement de $K$ dans $\dbC$ et $\psi = \exp(2\pi i x)$. Pour $\omega_{s}$, $dx$ la mesure de Lebesgue,
\[
\eps(\omega_{s},\psi,dx) = i^{N}.
\]
\end{para}

\begin{para}
\textbf{(3.4.2)} $K = \dbC$. Soient $z$ un plongement de $K$ dans $\dbC$ et $N \geq 0$. Pour $\psi = \exp(2\pi i \Tr_{\dbC/\dbR}(z))$ et $dx = |dz \wedge d\ov{z}| = 2\,da\,db$ (pour $z = a+bi$),
\[
\eps(z^{-N}\omega_{s},\psi,dx) = i^{N}.
\]
\end{para}

\begin{para}
\textbf{(3.4.3)} $K$ non archim\'edien. Si $\chi$ est non ramifi\'e, $dx = 1$ et $n(\psi) = 0$,
\begin{equation}\label{eq3.4.3.0}
\eps(\chi,\psi,dx) = 1.
\end{equation}
Pour $\chi$ ramifi\'e, on a
\begin{equation}\label{eq3.4.3.1}
\eps(\chi,\psi,dx) = \int_{y\scrO^{*}} \chi^{-1}(x) \psi(x)\, d^{*}x.
\end{equation}
On verra qu'il est naturel d'unifier ces deux formules de la fa\c{c}on suivante: si on d\'efinit $\eps_{0}(\chi,\psi,dx)$ par les formules
\begin{equation}\label{eq3.4.3.3}
\eps_{0}(\chi,\psi,dx) = \eps(\chi,\psi,dx) \quad \text{pour } \chi \text{ ramifi\'e, et } \eps_{0}(\chi,\psi,dx) = q^{n(\psi)} \quad \text{pour } \chi \text{ non ramifi\'e},
\end{equation}
on a
\begin{equation}\label{eq3.4.3.4}
\eps(\chi,\psi,dx) = \eps_{0}(\chi,\psi,dx) \cdot \chi(\pi^{n(\psi)+\sw(\chi)}) \cdot q^{-\sw(\chi)-n(\psi)} \cdot \frac{dx}{d^{*}x}.
\end{equation}
De (\ref{eq3.4.3.3}) et (\ref{eq3.4.3.4}), on d\'eduit que pour $\chi$ non ramifi\'e,
\begin{equation}\label{eq3.4.3.5}
\eps(\chi\omega_{s},\psi,dx) = \chi(\pi^{n(\psi)+a(\chi)}) \cdot q^{-(a(\chi)+n(\psi))\cdot s} \cdot \eps(\chi,\psi,dx).
\end{equation}
\end{para}

\begin{para}
Supposons $K$ non archim\'edien, et reprenons les notations \ref{para2.2.1}. Soit $V$ un espace vectoriel de dimension finie sur un corps $E$ de caract\'eristique $0$, ou sur $\dbC$. Les repr\'esentations $\rho: W(\ov{K}/K) \to \GL(V)$ seront toujours suppos\'ees continues (un sous-groupe ouvert de $I$ agit trivialement). Pour $V$ l'espace d'une repr\'esentation et $V_{I}$ le sous-espace des invariants sous l'inertie, on pose
\begin{equation}\label{eq3.5.1}
Z(V,t) = \det(1 - Ft, V_{I})^{-1}.
\end{equation}
On a
\begin{equation}\label{eq3.5.2}
L(V) = Z(V,1) \quad (E \subset \dbC \cup \{\infty\}).
\end{equation}
Notons $\omega_{t}: W(\ov{K}/K) \to E((t))^{*}$ le caract\`ere non ramifi\'e tel que $\omega_{t}(F) = t$ (en un sens \'evident, on a $\omega_{t} = \omega_{s}$ pour $t = q^{-s}$). On a
\begin{equation}\label{eq3.5.3}
Z(V,t) = L(V \otimes \omega_{t}).
\end{equation}
\end{para}

\begin{para}\label{para3.6}
C'est pour que (3.2.3) et (\ref{eq3.5.2}) soient compatibles que nous avons normalis\'e l'isomorphisme du corps de classe local de telle sorte que les Frobenius g\'eom\'etriques correspondent aux uniformisantes. Tant (3.2.3) que (\ref{eq3.5.2}) sont tr\`es naturels: (3.2.3) s'impose du point de vue analytique (cf.~(\ref{eq3.3.1})), et (\ref{eq3.5.2}) suit la tradition des g\'eom\`etres: c'est le Frobenius g\'eom\'etrique qui tend \`a agir sur la cohomologie des vari\'et\'es alg\'ebriques avec pour valeurs propres des entiers alg\'ebriques.
\end{para}

\begin{para}
Les repr\'esentations (continues) complexes irr\'eductibles de $W(\ov{K}/K)$ sont, pour $K$ complexe, les quasi-caract\`eres. Pour $K$ r\'eel, elles sont, outre les quasi-caract\`eres de $K^{*}$, les repr\'esentations induites par les quasi-caract\`eres $\chi$ de $\ov{K}^{*}$ tels que $\chi \not= \chi \circ \sigma$, pour $\sigma$ la conjugaison complexe de $\ov{K}$.
\end{para}

\begin{para}
Pour toute repr\'esentation complexe $V$, exprim\'ee, dans le groupe de Grothendieck des repr\'esentations de $W(\ov{K}/K)$ comme somme de repr\'esentations simples, on d\'efinit $L(V)$ comme suit:
\begin{enumerate}
\item[$1)$] $K$ complexe. Pour $V = [\chi]$ avec $\chi$ un quasi-caract\`ere de $\ov{K}^{*}$, $L(V) = L(\chi)$.
\item[$2)$] $K$ r\'eel. Outre les quasi-caract\`eres de $K^{*}$, il y a les repr\'esentations induites $\Ind_{\ov{K}^{*}}^{W}(\chi)$; on pose $L(V) = L(\chi)$.
\end{enumerate}
\end{para}

\begin{proposition}\label{prop3.8}
Soient $K$ un corps local et $\ov{K}$ une cl\^oture alg\'ebrique de $K$.
\begin{enumerate}
\item[\rm(i)] Pour toute suite exacte $0 \to V' \to V \to V'' \to 0$ de repr\'esentations complexes de $W(\ov{K}/K)$,
\begin{equation}\label{eq3.8.1}
L(V) = L(V') \cdot L(V'').
\end{equation}
\item[\rm(ii)] Soient $L$ une extension finie s\'eparable de $K$ dans $\ov{K}$ et $V_{L}$ une repr\'esentation complexe de $W(\ov{K}/L)$. Soit $V$ la repr\'esentation induite de $W(\ov{K}/K)$. On a
\begin{equation}\label{eq3.8.2}
Z(V,t) = Z(V_{L},t).
\end{equation}
\end{enumerate}
\end{proposition}

\begin{proof}
L'assertion (i) est triviale. Dans le cas archim\'edien, l'assertion (ii) r\'esulte de la formule de duplication
\[
\Gamma_{\dbC}(s) = \Gamma_{\dbR}(s) \cdot \Gamma_{\dbR}(s+1)
\]
correspondant \`a
\[
\Ind_{\ov{K}^{*}}^{W}([\omega_{s}]) = [\omega_{s}] + [\text{signe} \cdot \omega_{s}].
\]
Prouvons (ii) pour $K$ non archim\'edien. Notant $k$ et $\ell$ les corps r\'esiduels de $K$ et $L$, on a
\[
V_{I} = \Ind_{W(\ov{K}/L)}^{W(\ov{K}/K)}(V_{L})_{I} = \Ind_{W(\ov{k}/\ell)}^{W(\ov{k}/k)}((V_{L})_{I})
\]
(isomorphisme de $W(\ov{k}/k)$-repr\'esentations), les deux membres s'identifiant \`a
\[
H^{0}(I_{V_{L}})(\text{fonctions covariantes de } W(\ov{k}/k) \text{ dans } V_{L} \text{ agissant sur } W(\ov{k}/k) \text{ par translation}; \text{cf.~le diagramme commutatif})
\]
\[
\begin{CD}
W(\ov{K}/L) @>>> W(\ov{k}/\ell) \\
@VVV @VVV \\
W(\ov{K}/K) @>>> W(\ov{k}/k)
\end{CD}
\]
Ceci ram\`ene (\ref{eq3.8.2}) au lemme bien connu suivant.
\end{proof}

\begin{lemma}\label{lem3.9}
Soient $F$ un endomorphisme de $V$, $n$ un entier et $F_{n}$ l'endomorphisme de $V^{\otimes n} = V \otimes_{k[t]/(t^{n}-1)}$ tel que $F_{n}(v \otimes e_{i}) = v \otimes e_{i+1}$ ($0 \leq i < n-1$) et $F_{n}(v \otimes e_{n-1}) = F(v) \otimes e_{0}$. On a
\[
\det(1 - F_{n}t) = \det(1 - Ft^{n}).
\]
\end{lemma}

\begin{proof}
Par le principe de prolongement des identit\'es alg\'ebriques, on peut supposer $F$ diagonal, dans une base convenable de $V$; on se ram\`ene d\`es lors aussit\^ot au cas o\`u $\dim(V) = 1$, et o\`u $F$ est la multiplication par un nombre $a \not= 0$. Dans la base $f_{i} = a^{-i}(v \otimes e_{i})$, on a alors $F_{n}(f_{i}) = a f_{i+1}$ ($i \in \dbZ/n\dbZ$). Les vecteurs propres de $F_{n}$ sont les $(\sum_{i} \zeta^{-i} f_{i})$ ($\zeta$ caract\`ere de $\dbZ/n\dbZ$), et
\[
\det(1 - F_{n}t) = \prod_{\zeta^{n}=1} (1 - \zeta a t) = 1 - at^{n}. \qedhere
\]
\end{proof}

\begin{para}
Soit $K$ un corps global ($=$ $\dbA$-corps, au sens de \cite{15}). Pour toute place $v$ de $K$, on notera $K_{v}$ le corps local compl\'et\'e de $K$ en $v$. Les objets d\'efinis plus haut pour tout corps local seront, pour $K_{v}$, not\'es avec un indice $v$.
\end{para}

\begin{para}
On note $\dbA_{K}$ l'anneau des ad\`eles de $K$, $\vl x \vr$ la valeur absolue (v\'erifiant $\vl ax \vr = \vl a \vr \cdot \vl x \vr$ pour une mesure de Haar $dx$ sur $\dbA_{K}$), et $\omega_{s}$ le quasi-caract\`ere $\vl x \vr^{s}$ de $\dbA_{K}^{*}$. Notant $x_{v}$ la composante de $x \in \dbA_{K}$ dans $K_{v}$, on a donc
\[
\vl x \vr = \prod_{v} \vl x_{v} \vr_{v}.
\]
On d\'esignera par $dx$ l'unique mesure de Haar sur $\dbA_{K}$ telle que
\[
\int_{\dbA_{K}/K} dx = 1
\]
(mesure de Tamagawa), par $d^{*}x$ une mesure de Haar sur $\dbA_{K}^{*}$, et par $\psi$ un caract\`ere non trivial de $\dbA_{K}/K$, de composantes locales $\psi_{v}$. Pour $f \in C^{\infty}$ \`a support compact sur $\dbA_{K}$, on pose
\[
\wh{f}(y) = \int f(x) \psi(xy)\, dx.
\]
\end{para}

\begin{para}
L'\'equation fonctionnelle globale de Tate \cite{14} s'\'ecrit, pour $\chi$ un quasi-caract\`ere de $\dbA_{K}^{*}/K^{*}$,
\begin{equation}\label{eq3.11.1}
\int \wh{f}(x) \chi^{-1}(x) d^{*}x = \int f(x) \chi(x) d^{*}x
\end{equation}
(int\'egrales d\'efinies par prolongement analytique dans les familles $\chi \cdot \omega_{s}$). Posons
\begin{equation}\label{eq3.11.2}
L(\chi) = \prod_{v} L(\chi_{v}) \quad \text{(d\'efini par prolongement analytique)},
\end{equation}
et, pour $dx = \bigotimes_{v} dx_{v}$,
\begin{equation}\label{eq3.11.3}
\eps(\chi) = \prod_{v} \eps(\chi_{v}, \psi_{v}, dx_{v}).
\end{equation}
Le produit (\ref{eq3.11.3}) est ind\'ependant du choix de $\psi$ et de celui de la d\'ecomposition de $dx$, comme on d\'eduit de (\ref{eq3.3.2}) et (\ref{eq3.3.3}). Si, pour presque tout $v$, $dx_{v} = 1$, presque tous ses facteurs sont \'egaux \`a $1$.
\end{para}

Les formules (\ref{eq3.3.1}) et (\ref{eq3.11.1}) fournissent l'\'equation fonctionnelle globale des fonctions $L$ de Hecke
\begin{equation}\label{eq3.11.4}
L(\chi^{-1}) = \eps(\chi) \cdot L(\chi).
\end{equation}

\begin{para}
Les r\'esultats \'enonc\'es ci-dessous sont d\'emontr\'es dans \cite{1} et \cite{19}.
\end{para}

Soient $K$ un corps global, $\ov{K}$ une cl\^oture alg\'ebrique de $K$ et $V$ une repr\'esentation complexe de $\Gal(\ov{K}/K)$. (Toujours suppos\'ee continue: l'action de $\Gal(\ov{K}/K)$ se factorise par un groupe de Galois fini $\Gal(K'/K)$.)

Pour chaque place $v$ de $K$, soit encore $V_{v}$ une place de $\ov{K}$ au-dessus, et $V_{v}$ la repr\'esentation d\'eduite de $V$ par restriction au sous-groupe $W(\ov{K}_{v}/K_{v})$ du groupe de d\'ecomposition $\Gal(\ov{K}_{v}/K_{v})$.

\begin{enumerate}
\item[(A)] Pour $\chi$ un quasi-caract\`ere de $\dbA_{K}^{*}/K^{*}$, de composantes locales $\chi_{v}$, le produit infini
\[
L(V,\chi) = \prod_{v} L(V_{v} \otimes \chi_{v})
\]
peut \^etre d\'efini par prolongement analytique en $\chi$ dans les familles $\chi \cdot \omega_{s}$. Il converge pour $\Ree(s)$ grand, et se prolonge en une fonction m\'eromorphe de $s$.
\item[(B)] Notons $\check{V}$ le dual de $V$. On a
\begin{equation}
L(V,\chi) = \eps(V,\chi) \cdot L(\check{V},\chi^{-1}\omega_{1})
\end{equation}
avec $\eps(V,\chi) \in \dbC^{*}$ de la forme $A \cdot e^{Bs}$ pour $\chi$ variant dans les familles $\chi \cdot \omega_{s}$.
\item[(C)] Soit $L$ une extension finie de $K$ dans $\ov{K}$, $N_{L/K}$ la norme, $V_{L}$ une repr\'esentation de $\Gal(\ov{K}/L)$, et $V = \Ind_{\Gal(\ov{K}/L)}^{\Gal(\ov{K}/K)}(V_{L})$. Pour $\chi$ un quasi-caract\`ere de $\dbA_{K}^{*}/K^{*}$, on a
\[
\eps(V,\chi) = \eps(V_{L}, \chi \circ N_{L/K}).
\]
\end{enumerate}

(A) et (B) se d\'emontrent par r\'eduction au cas o\`u $\dim V = 1$, via le th\'eor\`eme de Brauer. (C) r\'esulte de (B) et de ce que (cons\'equence de \ref{lem3.9}),
\[
L(V,\chi) = L(V_{L}, \chi \circ N_{L/K}).
\]
On a aussi, de m\^eme,
\begin{equation}
\eps(V' + V'', \chi) = \eps(V',\chi) \cdot \eps(V'',\chi)
\end{equation}
ce qui permet de d\'efinir $\eps(V,\chi)$ pour $V$ seulement une repr\'esentation virtuelle.


\section{Existence des constantes locales}

\begin{theorem}\label{thm4.1}
Il existe une et une seule fonction $\eps$, v\'erifiant les conditions \textup{(1)} \`a \textup{(4)} ci-dessous, qui associe un nombre $\eps(V,\psi,dx) \in \dbC^{*}$ \`a chaque classe d'isomorphie de sextuples $(K, \ov{K}, \psi, dx, V, \rho)$ form\'es d'un corps local $K$, d'une cl\^oture alg\'ebrique $\ov{K}$ de $K$, d'un caract\`ere additif non trivial $\psi: K \to \dbC^{*}$, d'une mesure de Haar $dx$ sur $K$, d'un espace vectoriel $V$ de dimension finie sur $\dbC$ et d'une repr\'esentation $\rho: W(\ov{K}/K) \to \GL(V)$.

\textup{(1)} Pour toute suite exacte de repr\'esentations $0 \to V' \to V \to V'' \to 0$, on a
\[
\eps(V,\psi,dx) = \eps(V',\psi,dx) \cdot \eps(V'',\psi,dx).
\]
Cette condition montre que $\eps(V,\psi,dx)$ ne d\'epend que de la classe de $V$ dans le groupe de Grothendieck des repr\'esentations de $W(\ov{K}/K)$ et permet de donner un sens \`a $\eps(V,\psi,dx)$ pour $V$ seulement une repr\'esentation virtuelle.

\textup{(2)} $\eps(V,\psi,a\,dx) = a^{\dim V} \eps(V,\psi,dx)$. En particulier, pour $V$ virtuel de dimension $0$, $\eps(V,\psi,dx)$ est ind\'ependant de $dx$. On le note $\eps(V,\psi)$.

\textup{(3)} Si $L$ est une extension s\'eparable finie de $K$ dans $\ov{K}$ et que $V$ est la repr\'esentation virtuelle de $W(\ov{K}/K)$ induite par une repr\'esentation virtuelle de dimension z\'ero $V_{L}$ de $W(\ov{K}/L)$, on a
\[
\eps(V,\psi) = \eps(V_{L}, \psi \circ \Tr_{L/K}).
\]

\textup{(4)} Si $\dim V = 1$, $V$ \'etant d\'efini par un quasi-caract\`ere $\chi$ de $K^{*}$, on a $\eps(V,\psi,dx) = \eps(\chi,\psi,dx)$.
\end{theorem}

\begin{para}
Soit $v = \sum_{i,j} n_{ij}(H_{i},\chi_{ij}) \in R_{+}^{0}(W(\ov{K}/K))$ (\ref{variants1.10}). On suppose les $H_{i}$ deux \`a deux non conjugu\'es. Soit $K_{i}$ l'extension de $K$ d\'efinie par $H_{i}$, et notons encore $\chi_{ij}$ le quasi-caract\`ere de $K_{i}^{*}$ d\'efini par $\chi_{ij}$. Pour un choix quelconque de mesures de Haar sur les $K_{i}$, et $\psi$ un caract\`ere additif de $K$, on pose
\begin{equation}\label{eq4.2.1}
\eps(v,\psi) = \prod_{i,j} \eps(\chi_{ij}, \psi \circ \Tr_{K_{i}/K}, dx_{i})^{n_{ij}}.
\end{equation}
Il r\'esulte de (\ref{eq3.3.2}) que $\eps$ est ind\'ependant du choix des $dx_{i}$.
\end{para}

\begin{lemma}\label{lem4.3}
Avec la notation de l'\'enonc\'e, pour $v \in R_{+}^{0}(W(\ov{K}/K))$,
\[
\eps(v,\psi(ax)) = \det(v)(a) \cdot \eps(v,\psi).
\]
\end{lemma}

\begin{proof}
Il suffit de le v\'erifier pour $v$ de la forme $(H,\chi) - (H,1)$, $H$ correspondant \`a une extension $L$ de $K$ dans $\ov{K}$. Soit $v_{0}$ la repr\'esentation virtuelle $[\chi] - [1]$ de $W(\ov{K}/L)$. La formule (\ref{eq3.3.2}) fournit
\[
\eps(v,\psi(ax)) = \det(v)(a) \cdot \eps(v,\psi).
\]
On sait (\cite{10} XII \S4) que l'inclusion de $K^{*}$ dans $L^{*}$ s'identifie au transfert $W(\ov{K}/K)^{\ab} \to W(\ov{K}/L)^{\ab}$. Puisque $b(v) = \Ind(v_{0})$, \ref{lem4.3} r\'esulte de \ref{prop1.2}.
\end{proof}

\begin{lemma}\label{lem4.4}
Le th\'eor\`eme \ref{thm4.1} est vrai dans le cas archim\'edien.
\end{lemma}

\begin{proof}
Pour $K = \dbC$, on d\'efinit $\eps$ par les formules (1) et (4), et (2) est automatique. Pour $K = \dbR$, on v\'erifie sur \ref{para3.6} que l'application
\[
R_{+}^{0}(W(\ov{K}/K)) \longrightarrow R^{0}(W(\ov{K}/K))
\]
est surjective, de noyau engendr\'e par les $\chi_{s} - \chi_{t}$, pour
\[
((K^{*},\omega_{s}) - (K^{*},\omega_{t})) - (\sum_{r,s} \omega_{r} \cdot ([\chi\omega_{s}] - [\chi\omega_{s+1}]))
\]
Pour $V$ une repr\'esentation virtuelle de dimension $0$ de $W(\ov{K}/K)$ et $v \in R_{+}^{0}(W(\ov{K}/K))$ tel que $b(v) = V$, on pose
\begin{equation}\label{eq4.4.1}
\eps(V,\psi) = \eps(v,\psi).
\end{equation}
V\'erifions que le second membre est ind\'ependant du choix de $v$, i.e.~que pour $v$ v\'erifiant $b(v) = 0$, on a $\eps(v,\psi) = 1$. D'apr\`es \ref{lem4.3}, il suffit de le v\'erifier pour $\psi(x) = \exp(2\pi ix)$. Pour ce choix de $\psi$, l'assertion r\'esulte de la d\'etermination ci-dessus de $\Ker(b)$ et du fait que le second membre de (3.4.1), (3.4.2) soit ind\'ependant de $s$.

Il est clair que $\eps$ d\'efini par (\ref{eq4.4.1}) v\'erifie (1) \`a (4).
\end{proof}

On se restreint maintenant au cas non archim\'edien.

\begin{para}
Soit $K$ un corps local non archim\'edien, $\ov{K}$ une cl\^oture s\'eparable de $K$ et $K^{\nr} \subset \ov{K}$ l'extension non ramifi\'ee maximale de $K$. Soit $J$ un sous-groupe distingu\'e ouvert du groupe d'inertie $I = \Gal(\ov{K}/K^{\nr})$ et $L$ l'extension correspondante de $K^{\nr}$. Pour $\sigma \in I/J$, on pose
\[
t_{L/J}(\sigma) = \inf\{v'(\sigma(x) - x) \mid x \text{ entier dans } L\}
\]
et on d\'efinit des fonctions $a_{I/J}$ et $\sw_{I/J}$ par les formules
\begin{equation}\label{eq4.5.1}
a_{I/J}(\sigma) = -t_{L/J}(\sigma) \quad \text{pour } \sigma \not= e, \quad a_{I/J}(e) = \sum_{\sigma \not= e} t_{L/J}(\sigma)
\end{equation}
\[
\sw_{I/J} = a_{I/J} - (r_{I/J} - r_{I/J}^{I/J})
\]
pour $r_{I/J}$ la repr\'esentation r\'eguli\`ere.
\end{para}

D'apr\`es Artin (voir \cite{10} VI \S2 Th.~1), $a_{I/J}$ est le caract\`ere d'une repr\'esentation de $I/J$, la repr\'esentation d'Artin. D'apr\`es \textit{loc.~cit.}~prop.~2, $\sw_{I/J}$ est de m\^eme le caract\`ere d'une repr\'esentation $\Sw_{I/J}$, la repr\'esentation de Swan. Ces repr\'esentations ne sont d\'efinies qu'\`a isomorphisme pr\`es. $a_{I/J}$ est somme de $\Sw_{I/J}$ et de la repr\'esentation d'augmentation de $I/J$. Pour $J \subset J' \subset I$, on a (\textit{loc.~cit.}~prop.~3)
\begin{equation}\label{eq4.5.2}
a_{I/J}|_{I/J'} = a_{I/J'}, \quad \sw_{I/J}|_{I/J'} = \sw_{I/J'}.
\end{equation}
Si $d$ est la valuation du discriminant $\mathfrak{d}_{L/K}$, i.e.~la valuation de la diff\'erente $\mathfrak{D}_{L/K}$ (\cite{10} III \S3 prop.~6) et que $r$ d\'esigne une repr\'esentation r\'eguli\`ere, on a (\textit{loc.~cit.}~prop.~4)
\begin{equation}\label{eq4.5.3}
a_{I/J} = v_{K}(\mathfrak{d}_{L/K}) \cdot r_{I} + \Ind_{I/J}^{I}(a_{J})
\end{equation}
o\`u $a_{J}$ est la repr\'esentation d'augmentation de $J$.

\begin{para}
Soit $V$ une repr\'esentation de $W(\ov{K}/K)$, triviale sur $J \subset I$, et de caract\`ere $\chi$. On d\'efinit les conducteurs d'Artin et de Swan de $V$ par les formules
\begin{equation}\label{eq4.5.4}
a(V) = \dim \Hom_{I/J}(a_{I/J}, V) = \sw(V) + \dim V - \dim V_{I}
\end{equation}
\[
\sw(V) = \dim \Hom_{I/J}(\sw_{I/J}, V).
\]
D'apr\`es (\ref{eq4.5.2}), ces conducteurs sont ind\'ependants du choix de $J$.
\end{para}

Pour $V$ de dimension $1$, d\'efini par un quasi-caract\`ere $\chi$, on a (\textit{loc.~cit.}~prop.~5)
\begin{equation}\label{eq4.5.5}
a(V) = \text{conducteur } a(\chi) \text{ de } \chi.
\end{equation}
Les notations pr\'ec\'edentes sont donc compatibles avec celles introduites en \ref{sec3}.

\begin{lemma}\label{lem4.5}
Avec la notation de l'\'enonc\'e, pour $v \in R_{+}^{0}(W(\ov{K}/K))$ et $\chi$ un quasi-caract\`ere non ramifi\'e de $K^{*}$, on a
\[
\eps(v \cdot \chi, \psi, dx) = \chi(\pi)^{a(v) + \dim(v) \cdot n(\psi)} \cdot q^{-(a(v)+\dim(v) \cdot n(\psi))} \cdot \eps(v,\psi,dx).
\]
\end{lemma}

\begin{proof}
Il suffit de le v\'erifier pour $v$ de la forme $(H,\lambda) - (H,\mu)$, $H$ d\'efinissant une extension $L$ de $K$ et $\lambda, \mu$ \'etant deux quasi-caract\`eres de $L^{*}$. Soit $d$ la valuation de la diff\'erente de l'extension $L/K$, $f$ le degr\'e de l'extension r\'esiduelle, $I \subset W(\ov{K}/K)$ et $J \subset W(\ov{K}/L)$ les groupes d'inertie et $K' \subset J$ assez petit.

Calculons $a(v)$. La formule d'adjonction pour les repr\'esentations induites, (\ref{eq4.5.3}) et (\ref{eq4.5.5}) fournissent, pour tout quasi-caract\`ere $\nu$ de $H$ (i.e.~de $L^{*}$)
\[
a(\Ind_{H}^{G}([\nu])) = f \cdot a(\nu) + d \cdot \dim(\Ind_{H}^{G}([\nu]))
\]
d'o\`u
\[
a(v) = f(a(\lambda) - a(\mu)).
\]
Pour $\chi$ un quasi-caract\`ere de $L^{*}$, $\nu$ un quasi-caract\`ere non ramifi\'e de $L^{*}$, et $n = n(\psi \circ \Tr_{L/K})$ comme en \ref{sec3}, on a d'apr\`es (\ref{eq3.4.3.1}) ou (\ref{eq3.4.3.2}),
\[
\eps_{0}(\nu \cdot \chi, \psi \circ \Tr_{L/K}, dx) = \nu(\pi_{L})^{a(\chi)} \cdot \eps_{0}(\chi, \psi \circ \Tr_{L/K}, dx)
\]
d'o\`u
\[
\eps(v \cdot \chi, \psi, dx) = \chi(\pi)^{a(v)} \cdot \eps(v, \psi, dx)
\]
ce qui est la formule annonc\'e.
\end{proof}

\begin{terminology}
Soient $\ov{K}$ une extension galoisienne (finie ou non) de $K$, de groupe de Galois $G$, $\alpha$ un quasi-caract\`ere de $K^{*}$ et $\psi$ un caract\`ere additif non trivial de $K$. On dira qu'il existe une $\alpha,\psi$-th\'eorie des constantes pour $\ov{K}_{1}/K$ s'il existe $\eps_{\alpha,\psi}$ du type suivant:
\end{terminology}

\begin{enumerate}
\item[(0)] Pour $H$ un sous-groupe ouvert de $G$, d\'efinissant une extension finie $L$ de $K$, $V$ une repr\'esentation de $H$, et $dx$ une mesure de Haar sur $L$, $\eps_{\alpha,\psi}(V,dx) \in \dbC^{*}$ est d\'efini.
\item[(1)] Pour une suite exacte de repr\'esentations $0 \to V' \to V \to V'' \to 0$, $\eps_{\alpha,\psi}(V,dx) = \eps_{\alpha,\psi}(V',dx) \cdot \eps_{\alpha,\psi}(V'',dx)$. $\eps_{\alpha,\psi}$ garde donc un sens pour les repr\'esentations virtuelles.
\item[(2)] On a $\eps_{\alpha,\psi}(V,a\,dx) = a^{\dim V} \eps_{\alpha,\psi}(V,dx)$. En particulier, pour $V$ virtuel de dimension $0$, $\eps_{\alpha,\psi}$ est ind\'ependant de $dx$; on le note $\eps_{\alpha,\psi}(V)$.
\item[(3)] Pour $H_{1} \subset H_{2} \subset G$, d\'efinissant $L_{1} \supset L_{2}$, et $V$ une repr\'esentation virtuelle de dimension $0$ de $H_{1}$, on a
\[
\eps_{\alpha,\psi}(\Ind_{H_{1}}^{H_{2}}(V)) = \eps_{\alpha \circ N_{L_{2}/K}, \psi \circ \Tr_{L_{2}/L_{1}}}(V).
\]
\item[(4)] Pour $V$ de dimension $1$, d\'efini par un caract\`ere $\chi$ de $L^{*}$,
\[
\eps_{\alpha,\psi}(V,dx) = \eps(\alpha \circ N_{L/K} \cdot \chi, \psi \circ \Tr_{L/K}, dx).
\]
\end{enumerate}

\begin{lemma}\label{lem4.7}
Soient $\ov{K}_{1}/K$ et $\alpha$ comme plus haut, $\psi$ un caract\`ere additif non trivial de $K$. Il existe au plus une $\alpha,\psi$-th\'eorie des constantes pour $\ov{K}_{1}/K$. Pour qu'il en existe une, il faut et suffit que pour tout $v \in \Ker(R_{+}^{0}(G) \to R(G))$, on ait $\eps_{\alpha,\psi}(v) = 1$.
\end{lemma}

\begin{proof}
Soient $H \subset G$, $L$ et $V$ comme plus haut. D'apr\`es \ref{prop1.5}, il existe $v \in R_{+}^{0}(H)$ d'image $V - \dim(V) \cdot [1]$ dans $R(H)$. On a n\'ecessairement
\begin{equation}\label{eq4.7.1}
\eps_{\alpha,\psi}(V - \dim(V) \cdot [1]) = \eps_{\alpha,\psi}(v)
\end{equation}
d'o\`u l'unicit\'e. L'hypoth\`ese assure que $\eps_{\alpha,\psi}$, d\'efini par cette formule, est ind\'ependant du choix de $v$; on v\'erifie ais\'ement les axiomes (0) \`a (4).
\end{proof}

\begin{lemma}\label{lem4.8}
\begin{enumerate}
\item[\rm(i)] Soit $\alpha$ un quasi-caract\`ere non ramifi\'e de $K^{*}$ (resp.~$\psi' = \psi(ax)$ un caract\`ere additif non trivial). Pour qu'il existe une $\alpha,\psi$-th\'eorie des constantes pour $\ov{K}_{1}/K$, il faut et il suffit qu'il en existe une $\alpha\cdot\alpha',\psi$-th\'eorie (resp.~une $\alpha,\psi'$-th\'eorie). Pour $H$, $L$ et $V$ comme plus haut, $\psi_{L} = \psi \circ \Tr_{L/K}$, $\alpha_{L} = \alpha \circ N_{L/K}$ et $\alpha'_{L} = \alpha' \circ N_{L/K}$, on a alors
\begin{equation}\label{eq4.8.1}
\eps_{\alpha\cdot\alpha',\psi}(V,dx) = \alpha'(\pi)^{-\sw(V)} \cdot \det(V)(\pi)^{v(\alpha')} \cdot \eps_{\alpha,\psi}(V,dx)
\end{equation}
\begin{equation}\label{eq4.8.2}
\eps_{\alpha,\psi'}(V,dx) = \vl a \vr^{-\dim(V)} \cdot \det(V)(a) \cdot \eps_{\alpha,\psi}(V,dx).
\end{equation}
\end{enumerate}
\end{lemma}

\begin{proof}
La premi\`ere assertion r\'esulte du crit\`ere \ref{lem4.7} et de \ref{lem4.5} (resp.~de \ref{lem4.3}). Il suffit de prouver (\ref{eq4.8.1}) (\ref{eq4.8.2}) pour $K = L$. De plus, on peut supposer que soit $V \in R^{0}(\Gal(\ov{K}_{1}/K))$, soit $\dim V = 1$. Dans le premier cas, on applique \ref{lem4.5} et \ref{lem4.3}. Dans le second, on utilise (\ref{eq4.7.1}) et la d\'efinition \ref{eq3.4.3.3}.
\end{proof}

On dira qu'il existe une $\alpha$-th\'eorie des constantes pour $\ov{K}_{1}/K$ si pour tout (pour un) $\psi$, il existe une $\alpha,\psi$-th\'eorie des constantes pour $\ov{K}_{1}/K$.

\begin{reduction}\label{red4.9}
Le th\'eor\`eme \ref{thm4.1} est impliqu\'e par l'assertion suivante:
\begin{quote}
$(*)$ Pour tout corps local non archim\'edien $K$ et toute extension galoisienne finie $\ov{K}_{1}/K$, il existe un quasi-caract\`ere non ramifi\'e $\alpha$ de $K^{*}$, et une $\alpha$-th\'eorie des constantes pour $\ov{K}_{1}/K$.
\end{quote}
\end{reduction}

D'apr\`es \ref{lem4.8}, l'assertion $(*)$ implique que, pour $\ov{K}$ une cl\^oture s\'eparable de $K$ et $\alpha$ un quasi-caract\`ere non ramifi\'e quelconque de $K^{*}$, il existe une $\alpha$-th\'eorie des constantes pour $\ov{K}/K$.

\begin{para}
Toute repr\'esentation $V$ de $W(\ov{K}/K)$ est triviale sur un sous-groupe d'indice fini convenable $J$ de l'inertie $I$. Soit $F$ un Frobenius g\'eom\'etrique. Puisque $I/J$ est fini, une puissance convenable $F^{n}$ de $F$ agit trivialement sur $I/J$ par conjugaison, donc est centrale dans $W(\ov{K}/K)/J$. On dira que $V$ a un \emph{type} si une puissance $F^{m}$ de $F$ n'a qu'une valeur propre $a$. Le type de $V$ est alors l'\'el\'ement de la limite inductive
\[
\varinjlim_{n|m} \dbC^{*}, \quad \varphi_{n,m}: x \mapsto x^{m/n}
\]
d\'efini par $a \mapsto a^{m/n}$. Toute repr\'esentation irr\'eductible a un type. Les repr\'esentations de $W(\ov{K}/K)$ de type donn\'e $\tau$ forment une sous-cat\'egorie ab\'elienne $R_{\tau}$, et la cat\'egorie de toutes les repr\'esentations de $W(\ov{K}/K)$ est somme de ces sous-cat\'egories. Si $R_{\tau}(W(\ov{K}/K))$ est le groupe de Grothendieck de $R_{\tau}$, on a donc
\begin{equation}\label{eq4.10.1}
R(W(\ov{K}/K)) = \bigoplus_{\tau} R_{\tau}(W(\ov{K}/K)).
\end{equation}
Les repr\'esentations semi-simples de type $1$ sont exactement les repr\'esentations (continues) de $\Gal(\ov{K}/K)$. Si $\chi_{\tau}$ est un quasi-caract\`ere de type $\tau$ (il existe des quasi-caract\`eres, m\^eme non ramifi\'es, de tous les types), la torsion par $\chi_{\tau}$ d\'efinit donc un isomorphisme
\begin{equation}\label{eq4.10.2}
\cdot \chi_{\tau}: R_{1}(W(\ov{K}/K)) \isom R_{\tau}(W(\ov{K}/K)).
\end{equation}
\end{para}

\begin{para}
Soient $L$ une extension s\'eparable finie de $K$ dans $\ov{K}$ et $f$ le degr\'e de l'extension r\'esiduelle. Si une repr\'esentation $V$ de $W(\ov{K}/L)$ est de type $\sigma$, la repr\'esentation induite de $W(\ov{K}/K)$ est de type $\sigma^{f}$. On dit que $V$ est de $K$-type $\tau$ si $\sigma^{f} = \tau$. L'ensemble des repr\'esentations de $K$-type donn\'e est donc stable par induction.
\end{para}

\begin{variante}
Les arguments pr\'ec\'edents s'appliquent \`a tout groupe extension de $\hat{\dbZ}$ par un groupe fini, et aux limites projectives de telles extensions. Le cas des groupes de Weil globaux nous sera utile plus tard.
\end{variante}

\begin{para}
Admettons $(*)$. Pour tout type $\tau$, soit $\chi_{\tau}$ un caract\`ere non ramifi\'e de type $\tau$. Puisqu'il existe une $\chi_{\tau}^{-1}$-th\'eorie des constantes, on peut d\'efinir $\eps$ sur $R_{\tau}(W(\ov{K}/K))$ par la formule
\[
\eps(V,\psi,dx) = \eps_{\chi_{\tau}^{-1},\psi}(V \otimes \chi_{\tau}^{-1}, dx).
\]
Il r\'esulte de l'unicit\'e \ref{lem4.7} que cet $\eps$ est ind\'ependant du choix de $\chi_{\tau}$. On d\'efinit ensuite $\eps$ sur $R(W(\ov{K}/K))$ par additivit\'e (\ref{eq4.10.1}). D'apr\`es \ref{lem4.7}, s'il existe une th\'eorie \ref{thm4.1}, elle est donn\'ee par cet $\eps$.
\end{para}

\begin{para}
On a fait ce qu'il fallait pour que $\eps$ v\'erifie (1), (2), (4) et (3) lorsque $V_{L}$ a un type. Il reste \`a v\'erifier (3) pour $L$ de la forme $[\lambda \circ N_{L/K}] - [\mu \circ N_{L/K}]$, avec $\lambda$ et $\mu$ deux quasi-caract\`eres non ramifi\'es de $K^{*}$, de types donn\'es.

On a, par (\ref{eq3.4.3.1}), (\ref{eq3.3.2}), (\ref{eq3.3.3}),
\[
\eps(V,\psi) = \cdots
\]
D'apr\`es \ref{lem4.8}, on a d'autre part
\[
\eps_{\alpha,\psi}(V) = \cdots
\]
Soient $d$ la valuation de la diff\'erente de $L/K$, $e$ l'indice de ramification et $f$ le degr\'e de l'extension r\'esiduelle. La formule (\ref{eq4.5.1}) fournit
\[
a(v) = f \cdot d \cdot (\cdots)
\]
de sorte que la formule \`a d\'emontrer se r\'eduit \`a
\[
f \cdot d + [L:K] \cdot n(\psi) = 0
\]
qui est essentiellement la d\'efinition de la diff\'erente (\cite{10} III \S3).

Le point clef de la d\'emonstration de $(*)$ est le lemme suivant.
\end{para}

\begin{lemma}\label{lem4.12}
Soit $K'/K$ une extension galoisienne de corps globaux. Pour chaque place $w$ de $K$, on choisit une place de $K'$ au-dessus, et on la note encore $w$. Soient $\alpha$ un quasi-caract\`ere du groupe des classes d'id\`eles de $K$, et $v$ une place de $K$ qui ne se d\'ecompose pas dans $K'$. Si, pour toute place finie $w \not= v$ de $K$, il existe une $\alpha_{w}$-th\'eorie des constantes pour $\ov{K}_{w}/K_{w}$, alors, il existe une $\alpha_{v}$-th\'eorie des constantes pour $\ov{K}_{v}/K_{v}$.
\end{lemma}

\begin{proof}
Soient $\psi$ un caract\`ere non trivial de $\dbA_{K}/K$, et $\psi_{v}$ sa $v$-composante. Nous allons construire une $\alpha_{v},\psi_{v}$-th\'eorie des constantes pour $\ov{K}_{v}/K_{v}$.

On a $\Gal(K'/K)_{v} = \Gal(K'_{v}/K_{v})$. Un sous-groupe $H$ de $\Gal(K'/K)$ d\'efinit donc, outre une extension $L$ de $K$, une extension $L_{v}$ de $K_{v}$ de compl\'et\'e $L_{v}$.

Soit $dx$ la mesure de Tamagawa sur $\dbA_{L}$ (de volume $1$), et posons $dx = \bigotimes_{w} dx_{w}$ avec $dx_{v}$ choisi \`a l'avance. Toute repr\'esentation $V$ de $H$ d\'efinit une repr\'esentation $v$ de $\Gal(K'/L)$ et, par restriction, des repr\'esentations $V_{w}$ des groupes de d\'ecomposition $\Gal(K'_{w}/K_{w})$. Une constante globale $\eps(V,\alpha)$ a \'et\'e d\'efinie en \ref{sec3}. Soit $w$ une place de $L$, et notons encore $w$ la place de $K$ image. Pour $w$ infinie, on dispose (\ref{lem4.4}) de constantes locales $\eps(V_{w}, (\alpha \circ N_{L/K})_{w}, (\psi \circ \Tr_{L/K})_{w}, dx_{w})$. Pour $w$ finie autre que $v$, on dispose aussi par hypoth\`ese de constantes $\eps_{\alpha_{w},\psi_{w}}(V_{w}, dx_{w})$. Notant les unes et les autres $\eps_{\alpha_{w}}(V_{w},\psi_{w},dx_{w})$, on pose
\[
\eps_{\alpha_{v},\psi_{v}}(V,dx_{v}) = \frac{\eps_{\alpha}(V,\psi,dx)}{\prod_{w \not= v} \eps_{\alpha_{w}}(V_{w},\psi_{w},dx_{w})}.
\]
On v\'erifie que $\eps_{\alpha_{v},\psi_{v}}$ ne d\'epend pas du choix de la d\'ecomposition $dx = \bigotimes dx_{w}$ et v\'erifie les axiomes (0) \`a (4). Le point clef est 3.12(C).
\end{proof}

\begin{lemma}\label{lem4.13}
Soit $K'/K$ une extension galoisienne de corps locaux. Il existe une extension galoisienne de corps globaux $lK'/lK$, et une place $v$ de $lK$ qui ne se d\'ecompose pas dans $lK'$, telle que $K$ s'identifie au compl\'et\'e de $lK$ en $v$ et $K'$ \`a une sous-extension de $lK'_{v}/lK_{v}$.
\end{lemma}

\begin{proof}
C'est trivial: on prend d'abord $lK'/lK$ avec $lK_{v} \supset K' \supset K = lK_{v}$, puis on remplace $lK$ par l'extension d\'efinie par un groupe de d\'ecomposition en $v$.
\end{proof}

\begin{lemma}\label{lem4.14}
Soient $K$ un corps global, $S$ un ensemble fini de places finies de $K$ et $m: S \to \dbN$. Il existe un caract\`ere $\alpha$ du groupe des classes d'id\`eles tel que, pour $w \in S$, le conducteur de $\alpha_{w}$ soit $\geq m(w)$, et que pour toute place finie hors de $S$, $\alpha$ soit non ramifi\'e.
\end{lemma}

\begin{proof}
Si $K$ est un corps de nombres, on peut en effet librement sp\'ecifier $\alpha$ sur $\prod_{w \text{ fini}} \scrO_{w}^{*}$, car ce sous-groupe compact de $\dbA_{K}^{*}$ ne rencontre $K^{*}$ qu'en $\{1\}$.

Si $K$ est un corps de fonctions, de corps de constantes $k$, on a
\[
\prod_{w \text{ fini}} \scrO_{w}^{*} \cap K^{*} = k^{*}
\]
et on remarque qu'il existe des caract\`eres de $\prod_{w} \scrO_{w}^{*}$, de conducteur arbitrairement grand, qui soient triviaux sur $k^{*}$.
\end{proof}

D'apr\`es \ref{lem4.12}, \ref{lem4.13} et \ref{lem4.14}, l'assertion \ref{red4.9} $(*)$ et donc \ref{thm4.1} r\'esultent des deux lemmes suivants, o\`u $\ov{K}_{1}/K$ est une extension galoisienne finie de corps locaux non archim\'ediens, de groupe de Galois $G$ et $\alpha$ un quasi-caract\`ere de $K^{*}$.

\begin{lemma}\label{lem4.15}
Si $\alpha$ et $\ov{K}_{1}/K$ sont non ramifi\'es, il existe une $\alpha$-th\'eorie des constantes pour $\ov{K}_{1}/K$.
\end{lemma}

\begin{lemma}\label{lem4.16}
Pour $\ov{K}_{1}/K$ donn\'e, il existe $m$ tel que, si le conducteur $m$ de $\alpha$ est $\geq m$, il existe une $\alpha$-th\'eorie des constantes pour $\ov{K}_{1}/K$.
\end{lemma}

\begin{proof}[D\'emonstration de \ref{lem4.15}]
Choisissons pour $\psi$ un caract\`ere de $K$ tel que $n(\psi) = 0$. Il suffit alors de v\'erifier que pour $v \in \Ker(R_{+}^{0}(G) \to R(G))$, on a $\eps_{\alpha,\psi}(v) = 1$. Cela r\'esulte de ce que, pour $\chi$ non ramifi\'e, $\eps(\chi,\psi,dx) = 1$.
\end{proof}

\begin{proof}[D\'emonstration de \ref{lem4.16}]
Soit $\alpha$ de conducteur $m$ et $n = \left[\frac{m+1}{2}\right]$. La fonction $\alpha(1+a)$ est alors additive en $a$ pour $v(a) \geq m-n$. Si $\psi$ est un caract\`ere additif non trivial de $K$, il existe donc $y$, de valuation $-(m+n(\psi))$ et bien d\'efini modulo $\pi^{-n(\psi)}$, tel que
\[
\alpha(1+a) = \psi(ya) \quad \text{pour } v(a) \geq n.
\]
Pour $\chi$ un quasi-caract\`ere de $K^{*}$, de conducteur $\leq m-n$, on a alors
\begin{equation}\label{eq4.16.1}
\eps(\chi\alpha, \psi, dx) = \chi(y)^{-1} \cdot \eps(\alpha, \psi, dx).
\end{equation}
Pour $\pi^{m-n(\psi)} \not\subset y(\pi^{-(m-n)})$, l'int\'egrale partielle correspondante est nulle: c'est l'int\'egrale d'un caract\`ere additif non trivial. Puisque $\chi$ est constant de valeur $\chi(y)$ sur $y(\pi^{-(m-n)})$, on a donc
\[
\eps(\chi\alpha,\psi,dx) = \chi(y)^{-1} \cdot \eps(\alpha,\psi,dx).
\]

Nous allons montrer que, pour $m$ assez grand, on obtient une $\alpha,\psi$-th\'eorie des constantes pour $\ov{K}_{1}/K$ en posant, pour $H \subset G$ d\'efinissant une extension $L$ de $K$ et $V$ une repr\'esentation de $H$,
\[
\eps_{\alpha,\psi}(V,dx) = \det(V)(y)^{-1} \cdot \eps(\alpha \circ N_{L/K}, \psi \circ \Tr_{L/K}, dx)^{\dim V}.
\]
Les axiomes (0) (1) (2) sont triviaux. L'axiome (3) r\'esulte de ce que l'inclusion de $K^{*}$ dans $W(\ov{K}/L)^{\ab} \cong L^{*}$ s'identifie au transfert $W(\ov{K}/K)^{\ab} \to W(\ov{K}/L)^{\ab}$ (\cite{10} XIII \S4) et de \ref{prop1.2}. Reste l'axiome (4).

Soit $H \subset G$ correspondant \`a $L/K$, et $e$ l'indice de ramification. Rappelons que pour $u \in \scrO_{L}^{*}$ et $n = $ le plus petit entier tel que $n \geq 2v(u)$,
\begin{equation}\label{eq4.16.2}
N_{L/K}(1+u) = 1 + \Tr_{L/K}(u) \pmod{\pi^{n}}.
\end{equation}
On le v\'erifie en \'ecrivant la norme de $(1+u)$ comme le produit de ses conjugu\'es.

Nous noterons $O(1)$ tout nombre born\'e en terme de $\ov{K}_{1}/K$ seulement. D'apr\`es \ref{lem4.15.2}, si $\alpha$ est de conducteur $m$, $\alpha \circ N_{L/K}$ est de conducteur $e(m+O(1))$. De plus, pour $-y$ v\'erifiant $v(y) = -\frac{em}{2} + O(1)$,
\[
\alpha(1+\Tr_{L/K}(a)) = \psi(y \cdot \Tr_{L/K}(a))
\]
pour tout caract\`ere $\chi$ de $H$, de conducteur $O(1)$, on peut refaire le raisonnement fait plus haut; on obtient la formule voulue.
\end{proof}


\section{Formulaire}

\begin{para}
Outre la constante locale $\eps(V,\psi,dx)$ d\'efinie par \ref{thm4.1}, nous consid\'ererons aussi, dans le cas non archim\'edien, la constante $\eps_{0}$ d\'efinie par
\[
\eps_{0}(V,\psi,dx) = \eps(V,\psi,dx) \cdot \det(V)(\pi^{n(\psi)+\sw(V)}) \cdot q^{\sw(V)(n(\psi)+\sw(V)/2)} \cdot \left(\frac{dx}{d^{*}x}\right)^{\dim V}.
\]
\end{para}

Les constantes $\eps$ et $\eps_{0}$ ont les propri\'et\'es suivantes.

\begin{para}\label{para5.2}
$\eps$ et $\eps_{0}$ sont additifs en $V$, donc gardent un sens pour $V$ virtuel.
\end{para}

\begin{para}
$\eps$ et $\eps_{0}$ sont ind\'ependants de $dx$ pour $V$ virtuel de dimension $0$. De m\^eme pour $\eps$,
\begin{equation}
\eps(V,\psi(ax),dx) = \vl a \vr^{-\dim(V)} \det(V)(a)^{-1} \cdot \eps(V,\psi,dx)
\end{equation}
et de m\^eme pour $\eps_{0}$.
\end{para}

\begin{para}
Pour $K$ non archim\'edien et $\chi$ un quasi-caract\`ere non ramifi\'e de $K^{*}$, on a
\begin{equation}\label{eq5.5.1}
\eps(V \cdot \chi, \psi, dx, s) = \chi(\pi)^{a(V)+\dim(V) \cdot n(\psi)} \cdot q^{-(a(V)+\dim(V) \cdot n(\psi)) \cdot s} \cdot \eps(V,\psi,dx)
\end{equation}
et
\begin{equation}
\eps_{0}(V \cdot \chi, \psi, dx, s) = \chi(\pi)^{\sw(V)+\dim(V)(n(\psi)+1)} \cdot q^{-(\sw(V)+\dim(V)(n(\psi)+1)) \cdot s} \cdot \eps_{0}(V,\psi,dx).
\end{equation}
En particulier, si on pose
\[
\eps(V,\psi,dx,s) = \eps(V \cdot \omega_{s}, \psi, dx)
\]
et de m\^eme pour $\eps_{0}$, on a
\begin{equation}\label{eq5.5.2}
\eps(V,\psi,dx,s) = q^{-(a(V)+\dim(V) \cdot n(\psi)) \cdot s} \cdot \eps(V,\psi,dx)
\end{equation}
et
\[
\eps_{0}(V,\psi,dx,s) = q^{-(\sw(V)+\dim(V)(n(\psi)+1)) \cdot s} \cdot \eps_{0}(V,\psi,dx).
\]
Pour $\omega$ une repr\'esentation non ramifi\'ee, on a de m\^eme
\begin{equation}\label{eq5.5.3}
\eps(V \otimes \omega, \psi, dx) = \omega(\pi)^{a(V)} \cdot \det(V)(\pi)^{n(\psi)} \cdot q^{-a(V) \cdot s} \cdot \eps(V,\psi,dx).
\end{equation}
\end{para}

\begin{para}
Pour $L/K$ une extension s\'eparable finie, $V_{L}$ une repr\'esentation virtuelle de dimension $0$ de $W(\ov{K}/L)$ et $V$ la repr\'esentation induite de $W(\ov{K}/K)$, on a
\begin{equation}\label{eq5.6.1}
\eps(V,\psi) = \eps(V_{L}, \psi \circ \Tr_{L/K})
\end{equation}
et de m\^eme pour $\eps_{0}$. En d'autres termes, il existe $\lambda(L/K,\psi) \in \dbC^{*}$ tel que, pour $V$ de dimension quelconque,
\begin{equation}\label{eq5.6.2}
\eps(V,\psi) = \lambda(L/K,\psi)^{\dim(V)} \cdot \eps(V_{L}, \psi \circ \Tr_{L/K}).
\end{equation}
La m\^eme formule vaut pour $\eps_{0}$.

Dans Langlands \cite{9}, c'est (\ref{eq5.6.2}) qui fournit le cadre de la th\'eorie. De plus, il prend syst\'ematiquement pour $dx$ la mesure de Haar autoduale pour $\psi$, et il l'omet de la notation. Les constantes sont de plus normalis\'ees pour \^etre de valeur absolue complexe $1$. Tout ceci lui impose une autre formule (5.4).
\end{para}

\begin{para}
Soient $V^{*}$ le dual de $V$ et $dx'$ la mesure de Haar duale de $dx$, relativement \`a $\psi$. On a
\begin{equation}\label{eq5.7.1}
\eps(V,\psi,dx) \cdot \eps(V^{*}, \psi(-x), dx') = q^{a(V)+\dim(V) \cdot n(\psi)}.
\end{equation}
En particulier, si $K$ est non archim\'edien et que $V$ est unitaire, de sorte que $V^{*} = \ov{V}$, on a
\begin{equation}\label{eq5.7.2}
|\eps(V,\psi,dx)|^{2} = q^{a(V)+\dim(V) \cdot n(\psi)}.
\end{equation}
\end{para}

\begin{para}
$\eps$ appara\^it dans l'\'equation fonctionnelle locale de Tate; posant $\chi' = \omega_{1}/\chi$,
\begin{equation}\label{eq5.8.1}
\int \wh{f}(x) \chi'(x) d^{*}x = \eps(\chi,\psi,dx) \cdot \frac{\int f(x) \chi(x) d^{*}x}{L(\chi)} \quad \text{(Tate)}.
\end{equation}
Pour $K$ non archim\'edien, $y \in K^{*}$ de valuation $1+\sw(\chi)+n(\psi)$, on a
\begin{equation}\label{eq5.8.2}
\eps(\chi,\psi,dx)^{-1} = \chi(y) \int_{\scrO^{*}} \chi^{-1}(x) \psi(yx) d^{*}x.
\end{equation}
\end{para}

\begin{para}
Pour $K$ non archim\'edien, si $\chi$ est non ramifi\'e on a $\eps(\chi,\psi,dx) = \chi(\pi^{n(\psi)}) \cdot q^{n(\psi)} \cdot \int_{\scrO} \psi\, dx$. En particulier, si $n(\psi) = 0$ et $\int_{\scrO} dx = 1$, on a
\[
\eps(\chi,\psi,dx) = 1.
\]
\end{para}

\begin{para}
Prenons $K$ non archim\'edien, $a(\chi) = 1$, $n(\psi) = -1$ et $\int_{\scrO} dx = 1$. Soient $\chi_{k}$ et $\psi_{k}$ les caract\`eres de $k^{*}$ et $k$ d\'efinis par $\chi$ et $\psi$. On a
\begin{equation}\label{eq5.10.1}
\eps_{0}(\chi,\psi,dx) = -\tau(\chi_{k},\psi_{k})
\end{equation}
o\`u $\tau$ est la somme de Gauss
\begin{equation}\label{eq5.10.2}
\tau(\chi_{k},\psi_{k}) = -\sum_{x \in k^{*}} \chi_{k}(x)^{-1} \psi_{k}(x)
\end{equation}
(\'egale \`a $1$ pour $\chi$ trivial). En particulier,
\begin{equation}\label{eq5.10.3}
\eps(\chi,\psi,dx) = q^{-1} \cdot \tau(\chi_{k},\psi_{k}).
\end{equation}
\end{para}

\begin{para}
Soit $K$ un corps global. Pour $V$ une repr\'esentation du groupe de Weil global de $K$, $\psi$ un caract\`ere additif de non trivial de $\dbA_{K}/K$ et $dx = \bigotimes_{v} dx_{v}$ la mesure de Tamagawa sur $\dbA_{K}$, on pose
\begin{equation}
L(V) = \prod_{v} L(V_{v})
\end{equation}
et
\begin{equation}
\eps(V) = \prod_{v} \eps(V_{v}, \psi_{v}, dx_{v}).
\end{equation}
L'\'equation fonctionnelle globale s'\'ecrit
\begin{equation}
L(V^{*}) = \eps(V) \cdot L(V).
\end{equation}
\end{para}

\noindent
\textbf{D\'emonstrations:} \ref{para5.2}, 5.3, 5.6 et 5.8 sont \ref{thm4.1} (1) et (2), (3), (4) (et \ref{sec3}) --- 5.4 et 5.5 r\'esultent de \ref{lem4.8}.

Il suffit de prouver (\ref{eq5.7.1}) pour $V$ d\'efini par un caract\`ere; la formule r\'esulte alors de (\ref{eq5.8.1}) et de la formule d'inversion de Fourier. (\ref{eq5.7.2}) r\'esulte de (\ref{eq5.7.1}) et de (5.3) (5.4).

5.9 r\'esulte de (\ref{eq3.4.3.1}) et (\ref{eq5.10}) de (\ref{eq5.8.2}).

Enfin, 5.11 r\'esulte de ce qui pr\'ec\`ede et de (\ref{eq3.11.4}) par des arguments connus \cite{1}, \cite{19}.

La restriction de \ref{thm4.1} au cas des extensions et des caract\`eres mod\'er\'ement ramifi\'es \'equivaut aux identit\'es suivantes sur les sommes de Gauss (\ref{eq5.10.2}). On y d\'esigne par $k$ un corps fini \`a $q$ \'el\'ements et par $\psi$ un caract\`ere additif non trivial de $k$.

\begin{para}\textbf{(Hasse-Davenport).}
Soient $k'$ une extension de $k$ de degr\'e $n$ et $\chi$ un caract\`ere de $k^{*}$. On a
\begin{equation}
\tau(\chi \circ N_{k'/k}, \psi \circ \Tr_{k'/k}) = \tau(\chi,\psi)^{n}.
\end{equation}
\end{para}

\begin{para}
Soient $\chi$ un caract\`ere de $k^{*}$ et $n$ un entier premier \`a $q$. Soit $X$ un ensemble de repr\'esentants des classes d'isomorphie de couples $(k',\chi')$ form\'es d'une extension $k'$ de $k$ et d'un caract\`ere $\chi'$ de $k'^{*}$ tel que
\begin{enumerate}
\item[(a)] $\chi'|_{k^{*}} = \chi$,
\item[(b)] $\chi'$ n'est pas de la forme $\chi'' \circ N_{k'/k''}$ pour $k''$ interm\'ediaire entre $k$ et $k'$, i.e.~$\chi'^{q^{f'}} \not= \chi'$ pour $0 < f' < [k':k]$.
\end{enumerate}
Posons
\[
R(X,\psi) = \prod_{(k',\chi') \in X} \tau(\chi', \psi \circ \Tr_{k'/k}).
\]
Alors,
\begin{equation}
R(X,\psi)^{n} = \tau(\chi,\psi)^{|X|}.
\end{equation}
Cette identit\'e se d\'evisse en les deux suivantes.
\end{para}

\begin{para}\textbf{(Langlands \cite{9} 7.8).}
Soient $\ell$ un nombre premier \`a $q$, $f$ l'ordre de $q$ dans $(\dbZ/\ell)^{*}$, $k'$ une extension de $k$ de degr\'e $f$, $T$ un ensemble de repr\'esentants des orbites de $\Gal(k'/k)$ dans les caract\`eres (non triviaux) d'ordre $\ell$ de $k'^{*}$ et $\chi$ un caract\`ere de $k^{*}$. On a
\[
\prod_{\theta \in T} \tau(\chi \circ N_{k'/k} \cdot \theta, \psi \circ \Tr_{k'/k}) = \tau(\chi,\psi)^{|T|} \cdot \prod_{\theta \in T} \tau(\theta, \psi \circ \Tr_{k'/k}).
\]
\end{para}

\begin{para}\textbf{(Langlands \cite{9} 7.9).}
Soient $\ell$ un nombre premier divisant $q-1$, $T$ l'ensemble des caract\`eres d'ordre $\ell$ de $k^{*}$, $\chi$ un caract\`ere de $k^{*}$ qui n'est pas une puissance $\ell$-i\`eme, $k'$ une extension de degr\'e $\ell$ de $k$ et $\chi'$ un caract\`ere de $k'^{*}$ tel que $\chi'(x)^{\ell} = \chi \circ N_{k'/k}(x)$. On a
\[
\tau(\chi', \psi \circ \Tr_{k'/k})^{\ell} = \tau(\chi,\psi) \cdot \prod_{\mu \in T} \tau(\mu,\psi)^{-1}.
\]
\end{para}

\begin{para}\textbf{D\'emonstrations.}
Si $K'/K$ est une extension mod\'er\'ement ramifi\'ee de corps locaux non archim\'ediens et $\psi$ un caract\`ere additif de $K$ tel que $n(\psi) = -1$, on a encore $n(\psi \circ \Tr_{K'/K}) = -1$, de sorte que \ref{eq5.10} s'applique tant \`a $K$ qu'\`a $K'$.

Soit $V$ une repr\'esentation mod\'er\'ement ramifi\'ee de $W(\ov{K}/K)$. On peut toujours l'\'ecrire comme somme de repr\'esentations induites par des caract\`eres mod\'er\'es $\chi_{i}$ d'extensions non ramifi\'ees $K_{i}/K$:
\[
V = \dim(V) \cdot [1] + \sum_{i} \Ind([\chi_{i}] - [1]) + \sum_{i} (\Ind([1]) - [K_{i}:K][1]).
\]
Soient $k_{i}$ le corps r\'esiduel de $K_{i}$ et $\chi_{i}$ le caract\`ere de $k_{i}^{*}$ d\'efini par $\chi_{i}$. Pour $dx$ une mesure de Haar sur $K$ telle que $\int_{\scrO} dx = 1$, on trouve
\begin{equation}
\eps_{0}(V,\psi,dx) = (-1)^{\dim V} \cdot \prod_{i} \tau(\chi_{i} \circ \Tr_{k_{i}/k}, \psi_{k}).
\end{equation}
La formule (Hasse-Davenport) s'obtient en prenant une extension non ramifi\'ee $K'/K$, et en \'ecrivant que, pour $\chi$ un caract\`ere (ici mod\'er\'ement ramifi\'e) de $K^{*}$, on a
\[
\Ind_{K^{*}}^{K'^{*}}([\chi]) = \sum_{\mu} [\chi\mu] - [\chi]
\]
o\`u $\mu$ parcourt les caract\`eres non ramifi\'es d'ordre divisant $[K':K]$. Cette formule (Hasse-Davenport) assure que le second membre de (5.16.1) ne d\'epend que de $V$.

Les identit\'es 5.13 s'obtiennent en prenant une extension totalement ramifi\'ee de degr\'e $n$ $K'/K$, et en conjuguant (\ref{eq5.6.1}) et (5.16.1) pour $V$ induite par un caract\`ere mod\'er\'e $\chi$ de $K'$.

(5.14) et (5.15) sont le cas particulier de (5.13.1) o\`u $n$ est un nombre premier $\ell$. Dans 5.14 (resp.~5.15), on suppose que le caract\`ere $\chi$ de $K'^{*}$ est (resp.~n'est pas) une puissance $\ell$-i\`eme.
\end{para}


\section{R\'eduction modulo $\ell$ des constantes locales}

Soient $K$ un corps local non archim\'edien, $\ov{K}$ une cl\^oture s\'eparable de $K$ et reprenons les notations \ref{para2.2.1}. Soit $\Lambda$ un anneau commutatif dans lequel $p$ soit inversible. On simplifierait l'expos\'e, sans perte de g\'en\'eralit\'e, en supposant que $\Lambda$ est local, voire un corps. Pour $x \in K$, $\vl x \vr$ a par hypoth\`ese un sens dans $\Lambda$.

\begin{para}
Une mesure de Haar sur $K$, \`a valeurs dans un $\dbZ[1/p]$-module $M$ est une fonction simplement additive de parties compactes ouvertes de $K$, \`a valeurs dans $M$, et invariante par translation. La fonction additive d\'efinie par
\[
\mu_{0}(a + \pi^{n}\scrO) = \vl \pi^{n} \vr = q^{-n}
\]
est l'unique mesure de Haar \`a valeurs dans $\dbZ[1/p]$ telle que $\mu_{0}(\scrO) = 1$, et toute mesure de Haar \`a valeurs dans $M$ est de la forme $m \cdot \mu_{0}$ ($m = \mu_{0}(\scrO) \in M$). On peut int\'egrer une fonction localement constante \`a support compact \`a valeurs dans $\Lambda$ contre une mesure de Haar \`a valeurs dans $\Lambda$.
\end{para}

\begin{para}
Soient $\rho: W(\ov{K}/K) \to \GL_{\Lambda}(V)$ une repr\'esentation du groupe de Weil sur un $\Lambda$-module libre (ou projectif) $V$, et $J$ un sous-groupe distingu\'e ouvert de $I$ sur lequel $\rho$ soit trivial. Pour tout nombre premier $\ell \not= p$, on sait que la repr\'esentation de Swan $\Sw_{I/J}$ peut se r\'ealiser comme un $\dbZ_{\ell}[I/J]$-module projectif. Si $\Lambda$ est une $\dbZ_{\ell}$-alg\`ebre, $\Hom_{I/J}(\Sw_{I/J}, V)$ est donc un $\Lambda$-module projectif. Si $\Lambda$ est de plus local, il est libre; son rang est le conducteur de Swan de $V$. Plus g\'en\'eralement, on v\'erifie qu'on peut toujours d\'efinir le conducteur de Swan de $V$ comme une fonction localement constante \`a valeurs enti\`eres sur $\Spec(\Lambda)$. Ce conducteur est invariant par extension de l'anneau des scalaires $\Lambda$.
\end{para}

\begin{para}
Si $\zeta$ est racine $(p^{n})$-i\`eme de l'unit\'e dans $\Lambda$, l'ordre de $\zeta$ est une fonction localement constante sur $\Spec(\Lambda)$: deux quelconques des polyn\^omes cyclotomiques $\Phi_{p^{i}}(X)$ engendrent l'id\'eal trivial de $\Lambda[X]$ (puisque $p \in \Lambda^{*}$), de sorte que $\Lambda$ est produits d'anneaux $\Lambda_{i}$ ($1 \leq i \leq n$), la $i$\`eme coordonn\'ee de $\zeta$ v\'erifiant $\Phi_{p^{i}}(\zeta) = 0$. Sur $\Lambda_{i}$, $\zeta$ est d'ordre $p^{i}$.
\end{para}

\begin{para}
Si $\psi$ est un caract\`ere additif (continu) de $K$ \`a valeurs dans $\Lambda^{*}$, il est \`a valeurs dans les racines $(p^{n})$-i\`emes de l'unit\'e. Ceci permet de d\'efinir $n(\psi)$ comme une fonction localement constante sur $\Spec(\Lambda)$; en un id\'eal premier $\fht$, $n(\psi)$ est $+\infty$ si $\psi$ est trivial mod $\fht$, et le plus grand entier $n$ tel que $\psi|_{\pi^{-n}\scrO}$ soit trivial mod $\fht$ (ou dans le localis\'e $\Lambda_{\fht}$, c'est pareil). Par abus de langage, on dit que $n(\psi)$ est non trivial si $n(\psi)$ n'est nulle part $+\infty$.
\end{para}

\begin{para}
Pour $\chi \in \Hom(K^{*}, \Lambda^{*})$, $\psi$ un caract\`ere additif non trivial de $K$ \`a valeurs dans $\Lambda$, et $dx$ une mesure de Haar \`a valeurs dans $\Lambda$, on peut maintenant d\'efinir
\[
\eps_{0}(\chi,\psi,dx) \in \Lambda^{*}
\]
par la formule (\ref{eq3.4.3.4}) (plus pr\'ecis\'ement, si $\Spec(\Lambda)$ n'est pas connexe, on commence par \'ecrire $\Lambda = \prod \Lambda_{i}$, avec $n(\psi)$ et $\sw(\chi)$ constant sur $\Spec(\Lambda_{i})$. On d\'efinit alors la $i$\`eme projection de $\eps_{0}$ par (\ref{eq3.4.3.4})).
\end{para}

\begin{theorem}\label{thm6.5}
Il existe une et une seule fonction $\eps_{0}$ du type suivant:
\begin{enumerate}
\item[\rm(a)] Pour $K$ un corps local non archim\'edien, $\ov{K}$ une cl\^oture s\'eparable de $K$, $\Lambda$ un corps de caract\'eristique $\not= p$, $V$ une repr\'esentation de $W(\ov{K}/K)$ dans un vectoriel de dimension finie sur $\Lambda$, $\psi$ un caract\`ere additif non trivial \`a valeurs dans $\Lambda$ et $dx$ une mesure de Haar (non nulle) \`a valeurs dans $\Lambda$, on a $\eps_{0}(V,\psi,dx) \in \Lambda^{*}$.
\item[\rm(b)] $\eps_{0}$ est invariant par extension du corps $\Lambda$.
\item[\rm(c)] Soit $A$ un anneau de valuation discr\`ete de caract\'eristique r\'esiduelle $\ell$, $\eta$ son corps des fractions et $s = A/\mathfrak{m}$ son corps r\'esiduel. Pour $V$ une repr\'esentation de $W(\ov{K}/K)$ dans un $A$-module libre de rang fini, et pour $\psi$ et $dx$ \`a valeurs dans $A$, on a $\eps_{0}(V,\psi,dx) \in A^{*}$, et
\[
\eps_{0}(V_{s}, \psi_{s}, dx_{s}) \equiv \eps_{0}(V_{\eta}, \psi_{\eta}, dx_{\eta}) \pmod{\mathfrak{m}}.
\]
\item[\rm(d)] Les conditions (1) \`a (3) de \ref{thm4.1} sont v\'erifi\'ees.
\item[\rm(e)] Pour $V = [1]$, $\dim V = 1$, $\eps_{0}$ est donn\'e par 6.4.
\end{enumerate}
\end{theorem}

\begin{proof}
Par descente galoisienne, il suffit de traiter le cas o\`u $\Lambda$ est un corps s\'eparablement clos. Si $\Lambda$ est de caract\'eristique $0$, l'\'enonc\'e, \'etant purement alg\'ebrique, r\'esulte du cas $\Lambda = \dbC$ trait\'e en \ref{thm4.1}. Supposons donc $\Lambda$ de caract\'eristique $\ell \not= 0$.

Soient $G$ un quotient fini de $\Gal(\ov{K}/K)$ et $A$ un anneau de valuation discr\`ete d'in\'egale caract\'eristique de corps r\'esiduel $A/\mathfrak{m} = \Lambda$. On suppose $\psi$ relev\'e en un caract\`ere additif \`a valeurs dans $A^{*}$, $dx$ en une mesure de Haar \`a valeurs dans $A$, et que le corps des fractions $\eta$ de $A$ est assez gros relativement \`a $G$.

Pour tout caract\`ere $\chi$ d'une sous-groupe de $G$ \`a valeurs dans $\eta^{*}$, en fait dans $A^{*}$, la constante $\eps_{0}$ correspondante est clairement dans $A$. D'apr\`es (\ref{eq5.7.1}), elle est m\^eme dans $\scrO^{*}$. On en d\'eduit que, pour toute repr\'esentation de $G$ sur $\eta$, on a $\eps_{0}(V_{\eta}, \psi, dx) \in A^{*}$.

On sait que l'homomorphisme de d\'ecomposition $d: R_{+}^{0}(G) \to R_{+}^{\Lambda}(G)$ est surjectif (les repr\'esentations virtuelles se rel\`event). Pour $V$ et $\tilde{V} \in R_{+}^{0}(G)$ tel que $d(\tilde{V}) = V$, on pose
\[
\eps_{0}(V,\psi,dx) = \text{r\'eduction mod } \mathfrak{m} \text{ de } \eps_{0}(\tilde{V},\psi,dx).
\]
D'apr\`es \ref{prop1.8}, pour l\'egitimer cette d\'efinition, il suffit de v\'erifier que
\begin{quote}
$(*)$ Pour $H \subset G$ d\'efinissant une extension $L$ de $K$, $\psi_{L} = \psi \circ \Tr_{L/K}$ et $\chi', \chi''$ deux caract\`eres de $H$ \`a valeurs dans $\eta^{*}$, si $\chi' \equiv \chi''$ sont congrus mod $\mathfrak{m}$, on a
\[
\eps_{0}((H,\chi') - (H,\chi''), \psi_{L}) \equiv 1 \pmod{\mathfrak{m}}.
\]
\end{quote}
Les conducteurs de Swan de $\chi'$ et $\chi''$ sont tous deux \'egaux au conducteur de Swan de $\chi' \bmod \mathfrak{m}$. Pour $y = \pi^{1+\sw(\chi')+n(\psi_{L})}$, on a donc
\[
\eps_{0}(\chi',\psi_{L})^{-1} \cdot \eps_{0}(\chi'',\psi_{L}) = \frac{\int_{\scrO_{L}^{*}} \chi'^{-1}(x) \psi_{L}(yx)\, d^{*}x}{\int_{\scrO_{L}^{*}} \chi''^{-1}(x) \psi_{L}(yx)\, d^{*}x}.
\]
Num\'erateur et d\'enominateurs sont des unit\'es, et sont clairement congrus mod $\mathfrak{m}$. On a donc $\eps_{0} \equiv 1$.

Pour $V$ une repr\'esentation de $\Gal(\ov{K}/K)$, on obtient ainsi une fonction $\eps_{0}$ ayant les propri\'et\'es annonc\'ees. Ses propri\'et\'es formelles sont les m\^emes que celles explicit\'ees dans le formulaire 5.1 dans le cas $\Lambda = \dbC$. Pour passer de $\Gal(\ov{K}/K)$ \`a $W(\ov{K}/K)$ on peut donc reprendre l'argument utilis\'e en \ref{red4.9}, \`a l'aide de la deuxi\`eme formule (\ref{eq5.5.1}).
\end{proof}

\begin{remark}
Dans tout ce num\'ero, il a \'et\'e essentiel de travailler avec $\eps_{0}$ plut\^ot qu'avec $\eps$. D'abord pour avoir \`a utiliser le conducteur de Swan (invariant par r\'eduction mod $\ell$) plut\^ot que le conducteur d'Artin, ensuite pour n'avoir pas \`a distinguer, dans la d\'efinition de $\eps_{0}(\chi,\psi,dx)$, entre les cas non ramifi\'e et mod\'er\'ement ramifi\'e. La signification de $\eps_{0}$ appara\^itra plus clairement au num\'ero suivant.
\end{remark}


\section{Fonctions $L$ modulaires}

\begin{para}
Soit $X$ une courbe projective non singuli\`ere sur $\dbF_{p}$. On suppose $X$ irr\'eductible, on note $K$ son corps de fonctions rationnelles et on identifie points ferm\'es de $X$ et places de $K$. On ne suppose pas $X$ g\'eom\'etriquement irr\'eductible, i.e.~$\dbF_{p}$ alg\'ebriquement ferm\'e dans $K$.
\end{para}

\begin{para}
Soit $\Lambda$ un anneau noeth\'erien. On sait ce qu'est un faisceau constructible de $\Lambda$-modules sur $X$ (SGA 4 IX 2.3). Pour tout $x \in X$, et pour $\ov{k(x)}$ une extension s\'eparablement close de $k(x)$, un tel faisceau $\mathscr{G}$ a une fibre $\mathscr{G}_{\ov{x}}$ en $x$, qui est un $\Lambda$-module de type fini. Ce dernier ne d\'epend que de la cl\^oture s\'eparable de $k(x)$ dans $\ov{k(x)}$, et $\Gal(k(x)/k(x))$ agit sur $\mathscr{G}_{\ov{x}}$ par transport de structure. De plus, pour $v$ une place de $K$ (point ferm\'e de $X$), et $\ov{K}_{v}$ une cl\^oture alg\'ebrique de $K_{v}$, on dispose d'une fl\`eche de sp\'ecialisation
\[
\mathscr{G}_{\ov{k(v)}} \longrightarrow \mathscr{G}_{\ov{v}}
\]
Cette fl\`eche est \'equivariante; son image est donc contenue dans $\mathscr{G}_{\ov{v}}^{I_{v}}$. On dit que $\mathscr{G}$ est non ramifi\'e en $v$ si $\text{sp}_{v}$ est un isomorphisme. Un faisceau constructible est presque partout non ramifi\'e.
\end{para}

Une fois choisies une cl\^oture s\'eparable $\ov{K}$ de $K$ et une place de $\ov{K}$ au-dessus de chaque place de $K$, on peut r\'eciproquement d\'efinir un faisceau constructible de $\Lambda$-modules comme suit:
\begin{enumerate}
\item[(a)] On se donne les fibres $\mathscr{G}_{\ov{K}}$ et $\mathscr{G}_{\ov{k(v)}}$ (respectivement des $\Lambda$-$\Gal(\ov{K}/K)$-modules et $\Lambda$-$\Gal(\ov{k(v)}/k(v))$-modules, de type fini sur $\Lambda$).
\item[(b)] On se donne les fl\`eches de sp\'ecialisation (\'equivariantes pour l'action des groupes de d\'ecomposition)
\[
\text{sp}_{v}: \mathscr{G}_{\ov{k(v)}} \longrightarrow \mathscr{G}_{\ov{v}}.
\]
\end{enumerate}
Pour que (a) (b) d\'efinissent un faisceau, il suffit que les $\text{sp}_{v}$ soient presque tous des isomorphismes.

\begin{para}
Il nous sera commode de modifier la d\'efinition pr\'ec\'edente en rempla\c{c}ant les groupes de Galois par des groupes de Weil. Soit $Y$ un sch\'ema sur $\dbF_{p}$, $Y_{\text{\'et}}$ le site \'etale de $Y$, $S = \Spec(\dbF_{p})$ et $\ov{S} = \Spec(\ov{\dbF}_{p})$. Soient $\ov{\dbF}_{p}$ une cl\^oture alg\'ebrique de $\dbF_{p}$ et $\ov{Y} = Y \times_{\dbF_{p}} \ov{\dbF}_{p}$. Un faisceau de Weil (resp.~un faisceau de Weil de $\Lambda$-modules constructible) $\mathscr{G}$ sur $Y$ est un faisceau (resp.~un faisceau de $\Lambda$-modules constructible) $\mathscr{G}$ sur $\ov{Y}_{\text{\'et}}$ plus une action de Frobenius sur $\mathscr{G}$, au-dessus de son action sur $\ov{Y}$.

*Remarque (inutile): Les faisceaux de Weil en ensembles sur $Y$ forment un topos $Y_{\text{W\'et}}$. On peut le d\'ecrire comme un $2$-produit fibr\'e de topos: identifiant $\Spec(\ov{\dbF}_{p})_{\text{\'et}}$ au topos classifiant $B_{\hat{\dbZ}}$, on a
\[
Y_{\text{W\'et}} = Y_{\text{\'et}} \times_{B_{\hat{\dbZ}}} B_{\dbZ}.
\]*

Nous appellerons simplement faisceaux les faisceaux de Weil; les faisceaux sur $Y_{\text{\'et}}$ seront appel\'es faisceaux \'etales.
\end{para}

\begin{para}
Un faisceau de $\Lambda$-modules constructible peut se d\'ecrire en termes galoisiens comme en 7.2, en rempla\c{c}ant partout $\Gal(\ov{K}/K)$, $\Gal(\ov{K}_{v}/K_{v})$ et $\Gal(\ov{k(v)}/k(v))$ par $W(\ov{K}/K)$, $W(\ov{K}_{v}/K_{v})$ et $W(\ov{k(v)}/k(v)) \cong \dbZ$.
\end{para}

\begin{para}
Soit $K_{v}$ un corps local non archim\'edien. On d\'efinit un faisceau de Weil sur $\Spec(\scrO_{v})$ en paraphrasant 7.3. Soit $\ov{K}_{v}$ une cl\^oture s\'eparable de $K_{v}$.

Il revient au m\^eme de se donner un faisceau (de Weil) de $\Lambda$-modules constructible $\mathscr{G}$ sur $\Spec(\scrO_{v})$ ou de se donner
\begin{enumerate}
\item[(a)] la fibre $\mathscr{G}_{\ov{K}_{v}}$, un $\Lambda$-$W(\ov{K}_{v}/K_{v})$-module de type fini sur $\Lambda$;
\item[(b)] la fibre $\mathscr{G}_{\ov{k(v)}}$, un $\Lambda$-$W(\ov{k(v)}/k(v))$-module de type fini sur $\Lambda$;
\item[(c)] la fl\`eche de sp\'ecialisation, $W(\ov{K}_{v}/K_{v})$-\'equivariante
\[
\text{sp}: \mathscr{G}_{\ov{k(v)}} \longrightarrow \mathscr{G}_{\ov{K}_{v}}.
\]
\end{enumerate}
\end{para}

\begin{para}
Soient $v$ une place de $K$ et $x \in \Spec(\scrO_{v})$. Par restriction, tout faisceau (de Weil) de $\Lambda$-modules constructible sur $X$ en d\'efinit un sur $\Spec(\scrO_{v})$. Un faisceau est plat si ses fibres sont des $\Lambda$-modules plats.
\end{para}

\begin{para}
Soit $K_{v}$ un corps local non archim\'edien. On pose $\deg(v) = [k(v):\dbF_{p}]$ et on note $\omega_{t}$ le quasi-caract\`ere non ramifi\'e de $W(\ov{K}_{v}/K_{v})$, \`a valeurs dans $\Lambda(t)^{*}$, tel que $\omega_{t}(F_{v}) = t^{\deg(v)}$.
\end{para}

\begin{para}
Soit $\mathscr{G}_{v}$ un faisceau constructible plat sur $\Spec(\scrO_{v})$. On pose
\begin{equation}
Z(\mathscr{G}_{v},t) = \det(1 - F_{v}t^{\deg(v)}, \mathscr{G}_{\ov{k(v)}}) \cdot \det(1 - F_{v}t^{\deg(v)}, \mathscr{G}_{\ov{K}_{v}} \otimes \omega_{t})^{-1}
\end{equation}
et on d\'efinit un conducteur
\begin{equation}
a(\mathscr{G}_{v}) = \sw(\mathscr{G}_{\ov{K}_{v}}) + \dim(\mathscr{G}_{\ov{K}_{v}}) - \dim(\mathscr{G}_{\ov{k(v)}})
\end{equation}
(un entier si $\Lambda$ est local, une fonction localement constante sur $\Spec(\Lambda)$ en g\'en\'eral).
\end{para}

Supposons que $\Lambda$ soit un corps de caract\'eristique $\not= p$. Soient $\psi$ un caract\`ere additif non trivial de $K_{v}$ \`a valeurs dans $\Lambda^{*}$, et $dx$ une mesure de Haar \`a valeurs dans $\Lambda$, comme en \ref{thm6.5}. On pose
\begin{equation}
\eps(\mathscr{G}_{v},t) = \eps_{0}(\mathscr{G}_{\ov{K}_{v}} \otimes \omega_{t}, \psi, dx) \cdot \det(-F_{v}t^{\deg(v)}, \mathscr{G}_{\ov{K}_{v}})^{-1}.
\end{equation}
$\eps$ est un mon\^ome. D'apr\`es (\ref{eq5.5.2}), son degr\'e en $t$ est $a(\mathscr{G}_{v}) \cdot \deg(v)$.

\begin{remark}
Soit $M$ une repr\'esentation de $W(\ov{K}_{v}/K_{v})$. En termes galoisiens, on d\'efinit des faisceaux $j_{!}M$ et $j_{*}M$ sur $\Spec(\scrO_{v})$ par les fl\`eches de sp\'ecialisation
\[
\text{sp}: 0 \longrightarrow M \quad \text{et} \quad \text{sp}: M^{I_{v}} \longrightarrow M.
\]
Les quantit\'es $\eps$ et $\eps_{0}$ des \S\ pr\'ec\'edents s'identifient \`a la quantit\'e (7.6.3) relative respectivement \`a $j_{*}M$ et $j_{!}M$. Au contraire de $j_{*}$, le foncteur $j_{!}$ commute \`a la r\'eduction mod $\ell$. Ceci peut expliquer pourquoi, au \S6, nous n'avons pas d\'efini $\eps$.
\end{remark}

\begin{para}
Soit $\mathscr{G}$ un faisceau constructible de $\Lambda$-modules plats sur $X$. Pour chaque place $v$ de $K$, $\mathscr{G}$ induit un faisceau $\mathscr{G}_{v}$ sur $\Spec(\scrO_{v})$. On pose
\begin{equation}
Z(\mathscr{G},t) = \prod_{v} Z(\mathscr{G}_{v},t) \in \Lambda[[t]].
\end{equation}
\end{para}

\begin{para}
Supposons que $\Lambda$ soit un corps de caract\'eristique $\not= p$. Il existe sur l'anneau des ad\`eles de $K$ une unique mesure de Haar \`a valeurs dans $\dbZ[1/p]$ telle que $\int_{\dbA_{K}/K} dx = 1$. On peut \'ecrire $dx = \bigotimes_{v} dx_{v}$ ($dx_{v}$ = mesure de Haar sur $K_{v}$, \`a valeurs dans $\dbZ[1/p]$; pour presque tout $v$, $\int_{\scrO_{v}} dx_{v} = 1$). Les $dx_{v}$ d\'efinissent des mesures de Haar \`a valeurs dans $\Lambda$.

S'il existe un caract\`ere non trivial de $\dbA_{K}/K$, \`a valeurs dans $\Lambda^{*}$, on pose
\begin{equation}
\eps(\mathscr{G},t) = \prod_{v} \eps(\mathscr{G}_{v}, t).
\end{equation}
Presque tous les termes du produit sont \'egaux \`a $1$, et le produit ne d\'epend pas du choix de $\psi$ ni de la d\'ecomposition choisie de $dx$.

Soit $\Lambda'$ d\'eduit de $\Lambda$ par adjonction d'une racine primitive $p$-i\`eme de $1$. Il existe $\psi$ \`a valeurs dans $\Lambda'^{*}$. La constante globale correspondante, \'etant ind\'ependante de $\psi$, est invariante par $\Gal(\Lambda'/\Lambda)$, donc dans $\Lambda^{*}$. Ceci permet de d\'efinir $\eps$ m\^eme s'il n'existe pas de caract\`ere $\psi$.
\end{para}

\begin{para}
Notons $i$ l'inclusion de $\Spec(K)$ dans $X$. Si $V$ est une repr\'esentation de $W(\ov{K}/K)$, on note $i_{*}V$ le faisceau sur $X$ de fibre en $\ov{K}$ \'egale \`a $V$, et pour lequel les $\text{sp}_{v}$ sont $\text{sp}_{v}: V^{I_{v}} \hookrightarrow V$.
\end{para}

\begin{theorem}\label{thm7.11}
Soit $\Lambda$ un corps de caract\'eristique $\ell \not= p$. Soient $\mathscr{G}$ un faisceau constructible de $\Lambda$-vectoriels, $V$ la fibre de $\mathscr{G}$ en $\ov{K}$.
\begin{enumerate}
\item[\rm(i)] $Z(\mathscr{G},t)$ est une fraction rationnelle en $t$.
\item[\rm(ii)] Pour $t \to \infty$, $Z(\mathscr{G},t)$ est asymptotique \`a $\eps(\mathscr{G},t)$. En d'autres termes, $Z(\mathscr{G},t)/\eps(\mathscr{G},t)$ est une s\'erie formelle en $t^{-1}$, de terme constant $1$.
\item[\rm(iii)] On a l'\'equation fonctionnelle
\[
Z(\mathscr{G},t) = \eps(\mathscr{G},t) \cdot Z(\mathscr{G}^{*}, p^{-1}t^{-1}).
\]
\end{enumerate}
\end{theorem}

Soit $S$ un ensemble fini de places contenant toutes les places o\`u la repr\'esentation $V$ est ramifi\'ee. Soit $j$ l'inclusion de $X-S$ dans $X$, et soit $j_{*}V$ le faisceau sur $X$, de fibre g\'en\'erale $V$, non ramifi\'e en dehors de $S$, et de fibre nulle en $v \in S$. On v\'erifie aussit\^ot que (i) (ii) (iii) \'equivalent respectivement \`a:
\begin{enumerate}
\item[(i$'$)] $Z(j_{*}V,t)$ est fonction rationnelle de $t$;
\item[(ii$'$)] Pour $t \to \infty$, $Z(j_{*}V,t) \sim \eps(j_{*}V,t)$;
\item[(iii$'$)] $Z(j_{*}V,t) = \eps(j_{*}V,t) \cdot Z(j_{*}V^{*},p^{-1}t^{-1})$.
\end{enumerate}

Supposons que $\ell \not= 0$. Pour toute repr\'esentation $W$ de $W(\ov{K}/K)$, non ramifi\'e en dehors de $S$ (identifi\'ee \`a un faisceau localement constant sur $X-S$), posons
\begin{equation}
i_{*}W = \text{le faisceau nul en dehors de } S, \text{ de fibre en } v \in S \text{ \'egale } H^{1}(I_{v}, W).
\end{equation}
Soit $P' = \Ker(t_{\ell}: I_{v} \to \dbZ_{\ell}(1))$. Pour toute repr\'esentation $W$ de $I_{v}$ sur $\Lambda$, on a canoniquement
\begin{equation}
H^{1}(I_{v}, W) = H^{1}(P', W)^{\dbZ_{\ell}(1)} = (W_{P'}(-1))^{\dbZ_{\ell}(1)}(-1)
\end{equation}
(coinvariants tordu). Pour toute repr\'esentation $W$ de $\dbZ_{\ell}(1)$, on a canoniquement
\begin{equation}
H^{0}(\dbZ_{\ell}(1), W) = W^{\dbZ_{\ell}(1)} \quad \text{(invariants)}
\end{equation}
\begin{equation}
H^{1}(\dbZ_{\ell}(1), W) = W_{\dbZ_{\ell}(1)}(-1) \quad \text{(coinvariants tordu)}.
\end{equation}
\[
H^{i}(\dbZ_{\ell}(1), W) = 0 \quad (i \geq 2).
\]
Puisque la dualit\'e \'echange invariants et coinvariants, on a (cf.~2.3):

\begin{lemma}\label{lem7.11.4}
Pour $\ell \not= 0$, $H^{0}(I_{v}, W)$ et $H^{1}(I_{v}, W)$ sont canoniquement en dualit\'e.
\end{lemma}

Nous n'utiliserons pas que cette dualit\'e peut s'interpr\'eter comme un cup-produit \`a valeurs dans $H^{1}(I_{v}, \Lambda(1)) = \Hom(I_{v}, \Lambda)$, engendr\'e par $t_{\ell}$, mod $I_{v}$.

Pour $\ell = 0$, dans le courant de cette d\'emonstration, nous poserons
\[
H^{1}(I_{v}, W) = W_{P'}(-1)
\]
et d\'efinirons $Z_{v}(Rj_{*}W,t)$ par (7.11.1). Cette d\'efinition est justifi\'ee par 7.11.5 et 7.11.6 ci-dessous. Pour $W$ une repr\'esentation de $W(\ov{K}_{v}/K_{v})$, nous poserons encore
\[
Z_{v}(Rj_{*}W,t) = \det(1 - F_{v}t^{\deg(v)}, H^{0}(I_{v},W)) \cdot \det(1 - F_{v}t^{\deg(v)}, H^{1}(I_{v},W)(-1))^{-1}.
\]

\begin{lemma}\label{lem7.11.5}
Soient $A$ un anneau de valuation discr\`ete d'in\'egale caract\'eristique $\ell \not= p$ et $W$ un $A$-module libre sur lequel $W(\ov{K}_{v}/K_{v})$ agit. Soient $\eta$ le corps des fractions de $A$, $s = A/\mathfrak{m}$ le corps r\'esiduel. Alors
\[
Z_{v}(Rj_{*}W_{s},t) \equiv Z_{v}(Rj_{*}W_{\eta},t) \pmod{\mathfrak{m}}.
\]
\end{lemma}

\begin{proof}[Esquisse de d\'emonstration]
Soient $P' = \Ker(t_{\ell}: I_{v} \to \dbZ_{\ell}(1))$, $\sigma$ un g\'en\'erateur de $\dbZ_{\ell}(1)$, $F$ un Frobenius g\'eom\'etrique et $K$ le complexe concentr\'e en degr\'es $0$ et $1$
\[
K: \quad W \xrightarrow{\ \prod_{i=1}^{q-1} \sigma^{i}\ } W.
\]
L'endomorphisme $\sigma - 1$ de $W^{P'}$ est inversible: il suffit de le voir apr\`es r\'eduction modulo l'id\'eal maximal de $A$. Apr\`es r\'eduction, puisque $\sigma^{p^{n}} = 1$ est l'identit\'e, $\sigma - 1$ est unipotent et $q$ est l'unique valeur propre de $\sigma$ sur $W^{P'}$. 

Soit $(F_{0}, F_{1})$ l'endomorphisme de $K$ de composantes $\sigma$ et $q^{-1}\sigma$. La quantit\'e
\[
Z = \det(1 - Ft^{\deg(v)}, W^{P'})^{-1} \cdot \det(1 - F_{1}t^{\deg(v)}, W_{P'})
\]
est de formation compatible \`a l'extension des scalaires \`a $\eta$ ou $s$. Apr\`es extension des scalaires \`a $\eta$ (resp.~$s$), on a $Z = \det(1 - Ft^{\deg(v)}, H^{0}(K))^{-1} \cdot \det(1 - Ft^{\deg(v)}, H^{1}(K))$. On conclut en notant que
\[
H^{0}(K) = W^{I_{v}} \quad \text{et} \quad H^{1}(K) = W_{I_{v}}(-1).
\qedhere
\]
\end{proof}

\begin{variante}
Supposons $A$ complet et soit $W$ un $A$-module sur lequel $W(\ov{K}_{v}/K_{v})$ agit contin\^ument pour la topologie $\ell$-adique (et non plus, comme auparavant, pour la topologie discr\`ete). On pose encore
\[
H^{1}(I_{v}, W) = (W_{P'})^{\dbZ_{\ell}(1)}(-1) \quad \text{(coinvariants tordu)}
\]
et on d\'efinit $Z_{v}(Rj_{*}W,t)$ comme auparavant. Le $H^{1}$ d\'efini plus haut peut s'interpr\'eter comme un $H^{1}$ continu
\begin{equation}
H^{1}(I_{v}, W) = \varprojlim_{n} H^{1}(I_{v}, W/\ell^{n}W).
\end{equation}
La formule de r\'eduction \ref{lem7.11.5} reste valable (avec la m\^eme d\'emonstration).
\end{variante}

\begin{lemma}\label{lem7.11.7}
La formule (iii) \'equivaut \`a
\begin{equation}
Z(j_{*}V,t) = \eps(j_{*}V,t) \cdot Z(j_{*}V^{*}, p^{-1}t^{-1}).
\end{equation}
\end{lemma}

\begin{proof}
\ref{lem7.11.7} r\'esulte de l'identit\'e locale
\[
Z_{v}(j_{*}V,t) = \eps_{v}(V_{v} \otimes \omega_{t}, \psi_{v}, dx_{v}, t) \cdot Z_{v}(j_{*}V^{*}, p^{-1}t^{-1}).
\]
Si on remplace $V$ par la repr\'esentation $V_{v} \otimes \omega_{t}$ de $W(\ov{K}_{v}/K_{v})$ sur $\Lambda((t))$, cette formule se simplifie en
\[
\det(1 - F_{v}t^{\deg(v)}, V^{I_{v}}) \cdot \det(-F_{v}t^{\deg(v)}, V_{I_{v}}(-1)) = \eps_{0}(V_{v}, \psi_{v}, dx_{v}, t)
\]
qui r\'esulte de la dualit\'e locale \ref{lem7.11.4}.
\end{proof}

\begin{proof}[D\'emonstration de \ref{thm7.11}]
Prouvons \ref{thm7.11}. Si $\Lambda \subset \dbC$, le produit infini $Z(j_{*}V,t)$ converge pour $|t|$ assez petit, et
\[
Z(j_{*}V,t) = L(V,s) \quad (t = q^{-s})
\]
o\`u $L(V,s)$ est la fonction $L$ d'Artin-Weil. L'\'equation fonctionnelle montre que $Z$ est une fonction m\'eromorphe de $t$; l'\'equation fonctionnelle de la fonction $L$ fournit (iii).

Les \'enonc\'es (i) et (iii) sont purement alg\'ebriques. \'Etant vrais pour $\Lambda = \dbC$ (ainsi que (i$'$) et (7.11.8)), ils restent vrais pour $\Lambda$ de caract\'eristique $0$.

Supposons maintenant $\Lambda$ de caract\'eristique $\ell \not= 0$.

Soit $J$ un sous-groupe ouvert distingu\'e de $W(\ov{K}/K)$, contenant $I_{v}$ pour $v \in S$. Si une repr\'esentation $V$ de $W(\ov{K}/K)/J$ sur $\Lambda$ se rel\`eve en caract\'eristique $0$, les assertions (i$'$) et (7.11.8) pour $V$ s'obtiennent par r\'eduction mod $\ell$ (compte tenu, pour (7.11.8), de \ref{lem7.11.5}). Les quantit\'es $Z(j_{*}V,t)$, $Z(Rj_{*}V,t)$ et $\eps(j_{*}V,t)$ d\'ependent additivement de $V$ (une extension donne un produit). Pour obtenir (i) et (7.11.8) en g\'en\'eral il suffit donc de savoir que, apr\`es extension des scalaires toute repr\'esentation de $W(\ov{K}/K)/J$ se rel\`eve virtuellement en caract\'eristique $0$. La m\'ethode de \ref{red4.9} (cf.~4.10.3) ram\`ene cette question au cas des groupes finis, i.e.~au th\'eor\`eme de Brauer.

L'\'enonc\'e (ii$'$) r\'esulte de (7.11.8) et de ce que, pour $t \to 0$, $Z \to 1$.
\end{proof}

\begin{proposition}
Si $\Lambda$ est un corps de caract\'eristique $p$ et $\mathscr{G}$ un faisceau constructible de $\Lambda$-vectoriels, $Z(\mathscr{G},t)$ est encore une fraction rationnelle en $t$.
\end{proposition}

\begin{proof}
Dans la d\'emonstration de \ref{thm7.11} (i) nous n'avons en effet pas fait usage de l'hypoth\`ese $\ell \not= p$.
\end{proof}


\begin{thebibliography}{99}

\bibitem{1} E. Artin, \textit{Zur Theorie der $L$-Reihen mit allgemeinen Gruppencharakteren}, Abh. Math. Sem. Univ. Hamburg 8 (1931), 292--306 (={\OE}uvres, \textnumero 11).

\bibitem{2} M. Artin, A. Grothendieck, J.-L. Verdier, \textit{Th\'eorie des topos et cohomologie \'etale des sch\'emas} (SGA 4), Lecture Notes in Math. 269, 270, 305, Springer-Verlag 1972--1973.

\bibitem{3} P. Deligne, \textit{Formes modulaires et repr\'esentations $\ell$-adiques}, S\'em. Bourbaki 355 (1969), Lecture Notes in Math. 179, Springer-Verlag 1971.

\bibitem{4} P. Deligne, \textit{Formes modulaires et repr\'esentations de $\GL(2)$}, in this volume, p.~??.

\bibitem{5} P. Deligne, \textit{Les constantes des \'equations fonctionnelles}, notes manuscrites.

\bibitem{6} B. Dwork, \textit{On the Artin root number}, Amer. J. Math. 78 (1956), 444--472.

\bibitem{7} A. Grothendieck, \textit{Formule de Lefschetz et rationalit\'e des fonctions $L$}, S\'em. Bourbaki 279 (1965) (B.~S. M.~F., tome 91).

\bibitem{8} H. Jacquet and R.~P. Langlands, \textit{Automorphic forms on $\GL(2)$}, Lecture Notes in Math. 114, Springer-Verlag 1970.

\bibitem{9} R.~P. Langlands, \textit{On the functional equation of the Artin $L$-functions}, notes manuscrites.

\bibitem{10} J.-P. Serre, \textit{Corps locaux}, Hermann 1962.

\bibitem{11} J.-P. Serre, \textit{Repr\'esentations lin\'eaires des groupes finis}, Hermann 1967.

\bibitem{12} J.-P. Serre, J. Tate, \textit{Good reduction of abelian varieties}, Ann. of Math. 88 (1968), 492--517.

\bibitem{13} J. Tate, \textit{$p$-divisible groups}, in: Proc. of a Conference on Local Fields (Driebergen), Springer-Verlag 1967.

\bibitem{14} J. Tate, \textit{Fourier analysis in number fields and Hecke's zeta function}, thesis (1950), in: Algebraic Number Theory (Cassels-Frohlich), Academic Press 1967.

\bibitem{15} A. Weil, \textit{Basic number theory}, Springer-Verlag 1967.

\bibitem{16} A. Weil, \textit{Dirichlet series and automorphic forms}, Lecture Notes in Math. 189, Springer-Verlag 1971.

\bibitem{17} A. Weil, \textit{Sur les courbes alg\'ebriques et les vari\'et\'es qui s'en d\'eduisent}, Hermann 1948.

\bibitem{18} A. Weil, \textit{Numbers of solutions of equations in finite fields}, Bull. Amer. Math. Soc. 55 (1949), 497--508.

\bibitem{19} A. Weil, \textit{Sur la th\'eorie du corps de classe}, J. Math. Soc. Japan 3 (1951), 1--35.

\bibitem{20} P.~X. Gallagher, \textit{Determinants of representations of finite groups}, Abh. Math. Sem. Univ. Hamburg 28 (1965), 162--167.

\end{thebibliography}

\end{document}
